%----------------------------------------------------------------------------------------
%   Доорх хэсгийг өөрчлөх шаардлагагүй
%----------------------------------------------------------------------------------------
%!TEX TS-program = xelatex
%!TEX encoding = UTF-8 Unicode
\documentclass[12pt,A4]{report}

\usepackage{fontspec,xltxtra,xunicode}
\setmainfont[Ligatures=TeX]{Times New Roman}
\setsansfont{Arial}

% \usepackage[utf8x]{inputenc}
% \usepackage[mongohttps://www.overleaf.com/project/65b32d6efecd9d917f19312blian]{babel}
%\usepackage{natbib}
\usepackage{geometry}
%\usepackage{fancyheadings} fancyheadings is obsolete: replaced by fancyhdr. JL
\usepackage{fancyhdr}
\usepackage{float}
\usepackage{afterpage}
\usepackage{graphicx}
\usepackage{amsmath,amssymb,amsbsy}
\usepackage{dcolumn,array}
\usepackage{tocloft}
\usepackage{dics}
\usepackage{nomencl}
\usepackage{upgreek}
\newcommand{\argmin}{\arg\!\min}
\usepackage{mathtools}
\usepackage[hidelinks]{hyperref}

\usepackage{algorithm}
\usepackage{algpseudocode}
\usepackage{color, colortbl}
\definecolor{pink}{rgb}{0.97,0.59,0.5}

\usepackage{listings}
\DeclarePairedDelimiter\abs{\lvert}{\rvert}%
\makeatletter
\usepackage{caption}
\captionsetup[table]{belowskip=0.5pt}
\usepackage{subfiles}

\usepackage{listings}
\renewcommand{\lstlistingname}{Код}
\renewcommand{\lstlistlistingname}{\lstlistingname ын жагсаалт}

\usepackage{color}
\definecolor{codegreen}{rgb}{0,0.6,0}
\definecolor{codegray}{rgb}{0.5,0.5,0.5}
\definecolor{codepurple}{rgb}{0.58,0,0.82}
\definecolor{backcolour}{rgb}{0.99,0.99,0.99}
 
\lstdefinestyle{mystyle}{
    basicstyle=\ttfamily\small,
    backgroundcolor=\color{backcolour},   
    commentstyle=\color{codegreen},
    keywordstyle=\color{magenta},
    numberstyle=\tiny\color{codegray},
    stringstyle=\color{codepurple},
    %basicstyle=\footnotesize,
    breakatwhitespace=false,         
    breaklines=true,                 
    captionpos=b,                    
    keepspaces=false,                 
    numbers=left,                    
    numbersep=10pt,                  
    showspaces=false,                
    showstringspaces=true,
    showtabs=false,                  
    tabsize=2
}
 
\lstset{style=mystyle, label=DescriptiveLabel} 
\lstdefinelanguage{JavaScript}{
  keywords={break, case, catch, continue, debugger, default, delete, do, else, false, finally, for, function, if, in, instanceof, new, null, return, switch, this, throw, true, try, typeof, var, void, while, with},
  morecomment=[l]{//},
  morecomment=[s]{/*}{*/},
  morestring=[b]',
  morestring=[b]",
  ndkeywords={class, export, boolean, throw, implements, import, this},
  keywordstyle=\color{blue}\bfseries,
  ndkeywordstyle=\color{darkgray}\bfseries,
  identifierstyle=\color{black},
  commentstyle=\color{purple}\ttfamily,
  stringstyle=\color{red}\ttfamily,
  sensitive=true
}
\let\oldabs\abs
\def\abs{\@ifstar{\oldabs}{\oldabs*}}
\makenomenclature
\begin{document}


%----------------------------------------------------------------------------------------
%   Өөрийн мэдээллээ оруулах хэсэг
%----------------------------------------------------------------------------------------

% Дипломийн ажлын сэдэв
\title{Монгол Улсын хууль тогтоомжид үндэслэсэн, хиймэл оюунт, хууль зүйн зөвлөгөөний систем}
% Дипломын ажлын англи нэр
\titleEng{Artificial Intelligence-based Legal Advisory System Based on Mongolian Legislation}
% Өөрийн овог нэрийг бүтнээр нь бичнэ
\author{Гантулгын Жавхлан}
% Өөрийн овгийн эхний үсэг нэрээ бичнэ
\authorShort{Г. Жавхлан}
% Удирдагчийн зэрэг цол овгийн эхний үсэг нэр
\supervisor{О.Билгүүн}
% Хамтарсан удирдагчийн зэрэг цол овгийн эхний үсэг нэр
\cosupervisor{Др Ч.Алтангэрэл}

% СиСи дугаар 
\sisiId{22B1NUM5577}
% Их сургуулийн нэр
\university{МОНГОЛ УЛСЫН ИХ СУРГУУЛЬ}
% Бүрэлдэхүүн сургуулийн нэр
\faculty{МЭДЭЭЛЛИЙН ТЕХНОЛОГИ, ЭЛЕКТРОНИКИЙН СУРГУУЛЬ}
% Тэнхимийн нэр
\department{МЭДЭЭЛЭЛ, КОМПЬЮТЕРЫН УХААНЫ ТЭНХИМ}
% Зэргийн нэр
\degreeName{Дипломын ажлын тайлан}
% Суралцаж буй хөтөлбөрийн нэр
\programeName{Программ хангамж (D061303)}
% Хэвлэгдсэн газар
\cityName{Улаанбаатар}
% Хэвлэгдсэн огноо
\gradyear{2026 оны 03 сар}


%----------------------------------------------------------------------------------------
%   Доорх хэсгийг өөрчлөх шаардлагагүй
%----------------------------------------------------------------------------------------
\include{main-pre}

% Удиртгалыг оруулж ирэх ба abstract.tex файлд удиртгалаа бичнэ
\include{abstract}

%----------------------------------------------------------------------------------------
%   Дипломын үндсэн хэсэг эндээс эхэлнэ
%----------------------------------------------------------------------------------------
%\addcontentsline{toc}{part}{БҮЛГҮҮД}
% Шинэ бүлэг
\chapter{Үндэслэл}
\ \ \ \ \ \ \ \ \  \ \ \ \ Хиймэл оюун ухааны технологийн хөгжил сүүлийн
арван жилд эрчимжиж, ялангуяа том хэлний загвар (Large Language Model,
LLM)-уудын гарц мэдрэмжтэй, байгалийн хэлний текст үүсгэх чадварыг эрс
дээшлүүлсэн. OpenAI-н GPT цуврал, Google-н Gemini, Anthropic-н Claude зэрэг
загварууд асуултад хариулах, текст хураангуйлах, тайлбар гаргах зэрэг
олон төрлийн даалгаврыг хүний түвшинд гүйцэтгэх болжээ. Гэвч эдгээр
загваруудын нэг үндсэн сул тал нь мэдлэгийн огтлолын хязгаарлалт болон
нотлогдоогүй мэдээлэл буюу hallucination үүсгэх эрсдэл бөгөөд энэ нь
эрх зүйн мэдээлэл шиг нарийвчлал шаардсан чиглэлд онцгой аюул дагуулдаг.

\
\ \ \ \ \ \ \ \ \  \ \ \ \ Retrieval-Augmented Generation (RAG) технологи нь
энэ дутагдлыг арилгах зорилгоор 2020 оны үед Facebook-н судалгааны
лаборатори (Meta AI Research)-д анх санал болгогдсон бөгөөд өнөөдөр
мэдлэгт суурилсан системүүдийн суурь архитектур болж тогтжээ. RAG нь
хэрэглэгчийн асуултад хариулахаасаа өмнө эхлээд гадаад мэдлэгийн
сангаас хамааралтай баримтуудыг хайж олоод, тэдгээрийг LLM-д контекст
болгон өгснөөр эх сурвалжид тулгуурласан, баталгаажуулах боломжтой
хариулт гаргадаг. Энэ арга барил нь ялангуяа байнга өөрчлөгддөг, эсвэл
нарийн мэргэшсэн баримт бичгийн санд хандах шаардлагатай чиглэлүүдэд
өргөн хэрэглэгдэж байна.

\
\ \ \ \ \ \ \ \ \  \ \ \ \ Эрх зүйн салбарт RAG технологи хэрэглэх нь
дэлхийн хэмжээнд эрчимтэй хөгжиж байна. АНУ-д LexisNexis, Westlaw зэрэг
хуулийн мэдээллийн томоохон платформууд AI хайлтын системээ RAG-д
суурилсан болгон шинэчилсэн. Европт GDPR болон бусад хуулийн актын
тайлбарт зориулсан тусгай RAG системүүд хөгжүүлэгдсэн. Азид Хятад,
Япон, Солонгос зэрэг оронд ч үндэсний хэлний онцлогт тохирсон хуулийн
чатботууд идэвхтэй хэрэгжиж байна. Эдгээр туршлага нь RAG технологи
хуулийн мэдээллийн хүртээмжийг дээшлүүлэхэд үр дүнтэй гэдгийг
баталж байна.

\
\ \ \ \ \ \ \ \ \  \ \ \ \ Монгол улсын хувьд эрх зүйн мэдээллийн
цахим шилжилт чиглэлд тодорхой ахиц гарч, \texttt{legalinfo.mn} болон
\texttt{shuukh.mn} зэрэг нийтийн хандалттай мэдээллийн сангууд бий
болсон. Гэвч эдгээр эх сурвалж дахь мэдээллийг хиймэл оюун ухааны
аргаар боловсруулах, хэрэглэгчид ухаалгаар дамжуулах чиглэлд судалгаа
хийгдээгүй байна. Монгол хуулийн мэдээллийн сангийн хэмжээ жил ирэх
тусам нэмэгдэж байгаа бөгөөд 2024 оны байдлаар \texttt{shuukh.mn}-д
100,000 гаруй шүүхийн шийдвэр, \texttt{legalinfo.mn}-д 3,000 гаруй
хууль тогтоомжийн баримт бичиг байна. Энэ хэмжээний мэдээллийг хүн
гараар хайж судлах нь улам бүр боломжгүй болж байна.

\
\ \ \ \ \ \ \ \ \  \ \ \ \ Монгол хэлний онцлог нь RAG системийг
хэрэгжүүлэхэд тусгай бэрхшээл дагуулдаг. Монгол Кирилл үсгийн
Unicode боловсруулалтын асуудал, хуулийн нэр томьёоны стандарт бус
хэлбэр, зүйл заалтын дугаарлалтын онцлог, шүүхийн шийдвэр болон
хуулийн актын бүтцийн ялгаатай байдал зэрэг нь энгийн RAG системийг
шууд хэрэгжүүлэхэд хүрэлцэхгүй болгодог. Мөн Монгол хэлний
tokenizer болон embedding загвар хангалттай хөгжөөгүй байгаа нь
хайлтын чанарт сөргөөр нөлөөлж болзошгүй бөгөөд энэ асуудлыг
тусгайлан судлах шаардлагатай.

\
\ \ \ \ \ \ \ \ \  \ \ \ \ Иймд Монгол хуулийн чиглэлд RAG технологи
хэрэгжүүлэх нь зөвхөн техникийн бус, нийгмийн ач холбогдол бүхий
судалгаа болж байна. Хуулийн мэдээллийн хүртээмжийг нэмэгдүүлэх нь
эрх зүйн боловсролыг дэмжих, иргэдийг өөрийн эрхийнхээ талаар
мэдлэгтэй болгох, хуульчдын ажлын үр ашгийг дээшлүүлэх зэрэг олон
тийш эерэг нөлөө үзүүлнэ. Энэхүү судалгааны ажил нь тэр дундаа
Монгол хэлний NLP болон RAG технологийн уулзварт хийгдэх анхны
системчилсэн оролдлого болох учраас эрдэм шинжилгээний хувьд ч
шинэлэг хувь нэмэр болно гэж үзэж байна.

\chapter{Зорилго, зорилт}
\section{Зорилго}

Энэхүү дипломын ажлын үндсэн зорилго нь \texttt{shuukh.mn} болон 
\texttt{legalinfo.mn} эх сурвалжийн өгөгдөлд тулгуурлан хэрэглэгчид 
Монгол хэлээр хуулийн асуулт асууж, эх сурвалжид суурилсан хариулт болон 
лавлагаа авах боломжтой хиймэл оюун ухааны чатбот системийг боловсруулж, 
хэрэгжүүлж, үнэлэх явдал юм.

Дэд зорилгын хүрээнд дараах асуудлуудыг шийдвэрлэхийг зорив:

\begin{enumerate}
    \item Монгол хуулийн текстийн онцлогт тохирсон өгөгдөл цуглуулах, 
    цэвэрлэх, боловсруулах автомат систем боловсруулах.
    
    \item Хэрэглэгчийн асуултад хамгийн хамааралтай хуулийн мэдээллийг 
    олж авах хайлтын системийг хэрэгжүүлэх.
    
    \item Хариулт бүрт ашигласан хуулийн эх сурвалжийн лавлагааг 
    автоматаар бүрдүүлэх.
    
    \item Системийн хариултын үнэн зөв байдлыг туршилтаар шалгаж үнэлэх.
\end{enumerate}
\section{Зорилт}

Дэвшүүлсэн зорилгыг хангахын тулд дараах зорилтуудыг тавьж шийдвэрлэв:

\begin{enumerate}
  \item \textbf{Хуулийн мэдээлэл цуглуулах, боловсруулах системийг
        бүрдүүлэх.}
        \texttt{shuukh.mn} болон \texttt{legalinfo.mn} цахим хуудаснаас
        шүүхийн шийдвэр болон хуулийн актуудыг автоматаар татаж авах
        програм бичих. Татсан мэдээллийг цэвэрлэж, Монгол Кирилл
        үсгийн стандартчиллыг хангаж, системд ашиглахад бэлэн болгох.

  \item \textbf{Хуулийн мэдээллийг хайх системийг хэрэгжүүлэх.}
        Хэрэглэгчийн асуултын утга болон түлхүүр үгэд тулгуурлан
        хамааралтай хуулийн мэдээллийг хурдан, үнэн зөвөөр олж авах
        хайлтын систем боловсруулах. Олдсон үр дүнгүүдээс хамгийн
        тохиромжтой 5--8 хэсгийг сонгон авах шүүлтүүрийн арга
        нэмэх.

  \item \textbf{Хиймэл оюун ухаанд суурилсан хариулт үүсгэх системийг
        бүрдүүлэх.}
        Хайлтаар олдсон хуулийн мэдээлэлд тулгуурлан GPT-4o-mini
        хэлний загварыг ашиглан хэрэглэгчийн асуултад хариулах систем
        хэрэгжүүлэх. Систем нь зөвхөн эх сурвалжид байгаа мэдээлэлд
        үндэслэн хариулж, өөрөөс нь зохиож хариулахгүй байх нөхцөлийг
        хангах. Мөн хэрэглэгч нэг яриа дотор дараалсан олон асуулт
        асуух боломжтой байх.

  \item \textbf{Эх сурвалжийн лавлагаа харуулах арга боловсруулах.}
        Хариулт бүрт ашигласан хуулийн актын нэр, шүүхийн шийдвэрийн
        дугаар болон холбоосыг автоматаар гаргаж хэрэглэгчид харуулах
        систем хэрэгжүүлэх. Ингэснээр хэрэглэгч хариултын үнэн зөв
        эсэхийг эх сурвалжаас шалгах боломжтой болно.

  \item \textbf{Системийн хариултын чанарыг туршиж үнэлэх.}
        Монгол хуулийн чиглэлээр 50-аас доошгүй туршилтын асуулт
        бэлтгэж, системийн өгсөн хариулт болон лавлагааны үнэн зөв
        байдлыг шалгах. Үнэлгээний үр дүнд тулгуурлан системийн
        тохиргоог сайжруулж, эцсийн гүйцэтгэлийн шинжилгээ хийх.
\end{enumerate}

% \footnote{Datacare LLC: \url{https://datacare.mn/#services}}.
\newpage
\chapter{Хууль зүйн зөвлөгөө, эрх зүйн мэдээллийн системийн судалгаа}
\markboth{}{}
\thispagestyle{empty}

\section{Хиймэл оюун ухааныг хуульд хэрэглэсэн дэлхийн туршлага}

\ \ \ \ \ \ \ \ \  \ \ \ \ Хиймэл оюун ухааны технологи эрх зүйн салбарт
нэвтрэх үйл явц 2010-аад оны дундаас эрчимжиж, өнөөдөр дэлхийн олон
улс оронд хуулийн мэдээлэл хайх, гэрээ шинжлэх, шүүхийн шийдвэр
урьдчилан таамаглах зэрэг олон чиглэлд хиймэл оюун ухааны системүүд
идэвхтэй ашиглагдаж байна \cite{surden2019artificial}. Хуулийн
мэргэжилтнүүдийн ажлын ачааллыг бууруулах, иргэдэд эрх зүйн мэдээллийн
хүртээмжийг нэмэгдүүлэх зорилгоор хэрэгжүүлэгдэж буй эдгээр системүүд
нь технологийн хөгжлийн томоохон амжилт болж байна.

\ \ \ \ \ \ \ \ \  \ \ \ \ АНУ-д LexisNexis компани 2023 онд ``Lexis+~AI''
платформыг нэвтрүүлж, хуульчид асуулт асуухад холбогдох хуулийн актууд,
шүүхийн шийдвэрүүдийг эшлэлтэйгээр гаргаж өгдөг болгосон
\cite{lexisnexis2023}. Үүнтэй ижил чиглэлд Thomson Reuters компанийн
``Westlaw Precision'' систем нь GPT-4 загварыг хуулийн мэдээллийн
сантай хослуулан хуульч нарт хариулт болон эх сурвалжийн лавлагааг
нэгэн зэрэг гаргаж өгдөг \cite{thomsonreuters2023}. Эдгээр системүүд
нь хариулт бүрийг эх сурвалжид тулгуурлан баталгаажуулах боломжийг
хэрэглэгчид олгодог онцлогтой.

\ \ \ \ \ \ \ \ \  \ \ \ \ Европын холбоонд GDPR болон хуулийн актуудыг
тайлбарлах, хэрэгжүүлэхэд дэмжлэг үзүүлэх зорилгоор хэд хэдэн
мэргэшсэн систем хөгжүүлэгдсэн. Германд ``LegalBench'' болон ``JurBot''
системүүд хуулийн оюутнуудад болон залуу хуульчдад хуулийн тайлбар,
жишээ шийдвэрүүдийг хайж олоход туслалцаа үзүүлдэг
\cite{niklaus2021swiss}. Их Британид ``ROSS Intelligence'' систем нь
судалгааны асуултуудад хамааралтай хуулийн прецедент болон дүгнэлтүүдийг
хайж олж хуульчдад дэмжлэг үзүүлдэг байсан бөгөөд LexisNexis-т
нэгтгэгдсэн \cite{ross2023}.

\ \ \ \ \ \ \ \ \  \ \ \ \ Азийн орнуудад ч мөн адил хурдацтай хөгжил
ажиглагдаж байна. Хятадад ``Faxin'' болон ``Alpha'' системүүд хуулийн
мэргэжилтнүүдэд хэрэгт холбогдох хуулийн заалт, прецедентийг хайж олох,
шүүхийн шийдвэрийн хэв маягийг шинжлэхэд туслалцаа үзүүлдэг
\cite{he2021legal}. Японд ``LegalForce'' компани гэрээний баримт
бичгийг шинжлэх, эрсдэлийг илрүүлэх AI системийг амжилттай
хэрэгжүүлсэн бөгөөд өнөөдөр 3,000 гаруй хуулийн байгууллага ашигладаг
болжээ \cite{legalforce2022}. Солонгост ``Lbox'' платформ шүүхийн
шийдвэрийн мэдээллийн санд тулгуурлан хуульчдад хэрэг шийдвэрлэх
стратеги боловсруулахад дэмжлэг үзүүлдэг.

\ \ \ \ \ \ \ \ \  \ \ \ \ Иргэдэд чиглэсэн хуулийн мэдээллийн системүүдийн
чиглэлд ч томоохон дэвшил гарсан. Канадад ``DoNotPay'' аппликейшн нь
иргэдэд торгуулийг эсэргүүцэх, нислэгийн хугацаа хоцрогдсон тохиолдолд
нэхэмжлэл гаргах зэрэг энгийн хуулийн асуудлуудад автоматаар туслалцаа
үзүүлдэг \cite{donotpay2023}. Энэ мэт системүүд нь хуулийн зөвлөгөө
авахад шаардагдах зардал болон хугацааны хүндрэлийг мэдэгдэхүйц
бууруулж байна.

\section{Гадаадын ижил төстэй системүүдийн харьцуулалт}

\ \ \ \ \ \ \ \ \  \ \ \ \ Хуулийн чиглэлд хэрэгжиж буй хиймэл оюун
ухааны системүүдийг архитектур болон зорилгоор нь хэд хэдэн бүлэгт
ангилж болно. Хайлтад суурилсан системүүд, RAG архитектурт суурилсан
системүүд болон тусгай сургалттай хуулийн загварууд гэсэн гурван үндсэн
ангилал байдаг \cite{cui2023chatlaw}.

\ \ \ \ \ \ \ \ \  \ \ \ \ \ChatLaw нь Пекингийн их сургуулийн
судлаачдын боловсруулсан Хятадын хуулийн тусгай том хэлний загвар юм
\cite{cui2023chatlaw}. 300,000 гаруй хуулийн баримт бичгээр сургасан
энэхүү загвар нь Хятадын хуулийн онцлогт нийцсэн хариулт өгдөг.
Гэвч сургалтын өгөгдлийн огтлолын хязгаарлалт болон шинэ хуулийн
өөрчлөлтийг шуурхай тусгах чадваргүй байдал нь сул тал болдог.
RAG аргыг нэмэлтээр ашигласнаар энэ асуудлыг хэсэгчлэн шийдвэрлэсэн.

\ \ \ \ \ \ \ \ \  \ \ \ \ \LawGPT нь RAG архитектурт суурилсан
бөгөөд OpenAI-н загвартай хуулийн мэдээллийн санг хослуулсан систем юм
\cite{lawgpt2023}. Хэрэглэгчийн асуултыг векторт хөрвүүлж хуулийн
мэдээллийн сангаас хамааралтай хэсгүүдийг олсны дараа LLM-д дамжуулж
хариулт үүсгэдэг энэ систем нь эх сурвалжийн лавлагааг хариулттай хамт
харуулдаг онцлогтой. АНУ-ын хуулийн системд тохируулагдсан энэ систем
нь ойролцоо 85\%-ийн лавлагааны нарийвчлалтай ажилладаг болохыг
туршилт харуулсан.

\ \ \ \ \ \ \ \ \  \ \ \ \ \LEGAL-BERT нь хуулийн баримт
бичгүүдэд тусгайлан урьдчилсан сургалт хийсэн BERT загварын өргөтгөл
юм \cite{chalkidis2020legal}. Энгийн BERT-тэй харьцуулахад хуулийн
нэр томьёо, бүтцийг илүү сайн ойлгодог боловч ганцаарчлан хариулт
үүсгэх чадваргүй тул хайлтын системийн бүрэлдэхүүн болгон ашиглагдах
нь илүү тохиромжтой.

\ \ \ \ \ \ \ \ \  \ \ \ \ Дараах хүснэгтэд дэлхийн тэргүүлэх хуулийн
AI системүүдийн үндсэн шинж чанарыг харьцуулсан болно:

\begin{table}[h]
\centering
\caption{Хуулийн AI системүүдийн харьцуулалт}
\begin{tabular}{|l|l|l|l|l|}
\hline
\textbf{Систем} & \textbf{Архитектур} & \textbf{Хэл} &
\textbf{Лавлагаа} & \textbf{Нээлттэй} \\
\hline
LexisNexis AI  & RAG + LLM   & Англи   & Тийм         & Үгүй \\
Westlaw        & RAG + GPT-4 & Англи   & Тийм         & Үгүй \\
ChatLaw        & Fine-tuned  & Хятад   & Хэсэгчлэн    & Тийм \\
LawGPT         & RAG         & Англи   & Тийм         & Хэсэгчлэн \\
LEGAL-BERT     & Embedding   & Олон    & Үгүй         & Тийм \\
Huuli.tech     & RAG         & Монгол  & Тийм         & Үгүй \\
\hline
\end{tabular}
\label{tab:legal_ai_comparison}
\end{table}

\section{Монголын эрх зүйн мэдээллийн одоогийн байдал}

\ \ \ \ \ \ \ \ \  \ \ \ \ Монгол улсад эрх зүйн мэдээллийн цахим шилжилт
чиглэлд сүүлийн жилүүдэд тодорхой ахиц дэвшил гарсан. 2000-аад оноос
эхлэн хуулийн баримт бичгүүдийг цахимжуулах ажил эхэлж, өнөөдөр
хэд хэдэн нийтийн хандалттай мэдээллийн сан бий болжээ
\cite{monlegalinfo2020}.

\ \ \ \ \ \ \ \ \  \ \ \ \ \legalinfo цахим хуудас нь Монгол
улсын Хууль зүй, дотоод хэргийн яамны санаачилгаар хөгжүүлэгдсэн бөгөөд
хуулийн акт, тогтоол, шийдвэр, дүрэм журмын мэдээллийн санг нийтэд
нээлттэй байдлаар санал болгодог. 2024 оны байдлаар 3,000 гаруй хуулийн
баримт бичиг, 15,000 гаруй дүрэм журам энд байршиж байна. Гэвч уг
системийн хайлтын чадавхи хязгаарлагдмал бөгөөд хэрэглэгч яг таарсан
нэр томьёогоор хайх шаардлагатай байдаг. Утга санааны дагуу хайх буюу
``би гэр орноо хэрхэн хамгаалах вэ'' гэх мэт байгалийн хэлний асуултад
хариулах чадваргүй хэвээр байна.

\ \ \ \ \ \ \ \ \  \ \ \ \ \shuukh.mn цахим хуудас нь Монгол улсын
Дээд шүүхийн санаачилгаар бий болсон шүүхийн шийдвэрийн мэдээллийн
сан юм. 2024 оны байдлаар 100,000 гаруй шүүхийн шийдвэр байршиж байгаа
бөгөөд иргэний хэрэг, эрүүгийн хэрэг, захиргааны маргаан гэсэн ангиллаар
хайх боломжтой. Гэвч шийдвэрийн бүтэц, форматыг стандартчилаагүй,
хайлтын боломж хязгаарлагдмал байгаа нь ижил төстэй прецедент хайх
болон шинжлэхэд хүндрэл учруулж байна.

\ \ \ \ \ \ \ \ \  \ \ \ \ Эдгээр эх сурвалжийн нийтлэг сул тал бол
хэрэглэгч хайж буй зүйлээ яг таарсан нэр томьёогоор мэддэг байх
шаардлагатай байдаг явдал юм. Монгол хэлний хайлтын системүүд
морфологийн нарийн өөрчлөлтийг ихэнх тохиолдолд зөв боловсруулж
чаддаггүй бөгөөд энэ нь хайлтын үр дүнгийн чанарт сөргөөр нөлөөлдөг
\cite{purev2008mongolian}. Тухайлбал, ``гэрээ'' болон ``гэрээний'',
``хариуцлага'' болон ``хариуцлагатай'' гэсэн хэлбэрүүдийг ижил утгатай
гэж таних чадвар дутмаг байдаг нь хайлтын үр нөлөөг бууруулдаг.

\section{Huuli.tech — Монголын хуулийн AI системийн туршлага}

\ \ \ \ \ \ \ \ \  \ \ \ \ \Huuli.tech нь Монгол улсад хуулийн
чиглэлд хиймэл оюун ухааны технологийг хэрэгжүүлсэн анхдагч системүүдийн
нэг бөгөөд \texttt{huuli.tech/chat} хаягаар хэрэглэгчдэд нээлттэй байдаг.
Энэхүү систем нь Монгол хуулийн мэдээллийн санд тулгуурлан хэрэглэгчийн
асуултад хариулах боломжийг олгодог бөгөөд Монголд энэ чиглэлийн
системийг хэрэгжүүлэх боломжтойг анх нотолсон туршлага болжээ.

\ \ \ \ \ \ \ \ \  \ \ \ \ Huuli.tech системийн үндсэн онцлогт дараах
зүйлс орно. Нэгдүгээрт, хэрэглэгч Монгол хэлээр байгалийн хэлний
асуулт бичиж болдог бөгөөд систем холбогдох хуулийн заалтыг хайж олж
хариулт гаргана. Хоёрдугаарт, хариулт бүрд ашигласан хуулийн эх
сурвалжийн лавлагаа харуулдаг тул хэрэглэгч хариултыг баталгаажуулах
боломжтой. Гуравдугаарт, системийн хариулт нь \texttt{legalinfo.mn}
болон бусад хуулийн эх сурвалжид суурилсан тул мэдээллийн үнэн зөв
байдлыг хангахыг зорьсон байдаг.

\ \ \ \ \ \ \ \ \  \ \ \ \ Гэвч Huuli.tech системтэй харьцуулахад
энэхүү дипломын ажлын хүрээнд боловсруулж буй систем хэд хэдэн
чухал ялгаатай. Нэгдүгээрт, \texttt{shuukh.mn} дахь шүүхийн шийдвэрийн
өгөгдлийг нэмж оруулснаар хуулийн акт болон шүүхийн практикийг хослуулан
хариулт гаргах боломжийг нэмнэ. Хоёрдугаарт, утга санааны хайлт болон
түлхүүр үгийн хайлтыг хослуулсан аргыг хэрэгжүүлснээр хайлтын чанарыг
дээшлүүлнэ. Гуравдугаарт, олон эргэлтийн яриа болон ярианы түүхийг
хадгалах аргаыг нэмснээр хэрэглэгчийн туршлагыг сайжруулна.
Дөрөвдүгээрт, системийн архитектурыг бүрэн нээлттэй, баримтжуулсан
байдлаар боловсруулж цаашид судалгааны үндэс болох боломжтой болгоно.

\section{RAG технологийн онолын үндэс}

\ \ \ \ \ \ \ \ \  \ \ \ \ Retrieval-Augmented Generation (RAG) технологийг
анх 2020 онд Meta AI Research-ийн судлаачид Patrick Lewis болон
бусад хамтрагчид ``Retrieval-Augmented Generation for Knowledge-Intensive
NLP Tasks'' нэртэй судалгааны бүтээлдээ танилцуулсан
\cite{lewis2020retrieval}. Уг аргын үндсэн санаа нь том хэлний загварын
(LLM) текст үүсгэх чадварыг баримт бичгийн санд суурилсан мэдээлэл хайх
системтэй нэгтгэх явдал юм. Хэрэглэгчийн асуулт ирэхэд эхлээд мэдлэгийн
сангаас хамааралтай мэдээллийг хайж олоод, тэдгээрийг LLM-д нэмэлт
контекст болгон өгснөөр нотолгоотой, эх сурвалжид тулгуурласан хариулт
гаргадаг.

\ \ \ \ \ \ \ \ \  \ \ \ \ RAG системийн архитектур нь дараах гурван үндсэн
бүрэлдэхүүнтэй \cite{gao2023retrieval}. Нэгдүгээрт,
өгөгдөл боловсруулах хэсэг нь эх сурвалжийн баримт бичгүүдийг
жижиг хэсгүүдэд хуваах, векторт хөрвүүлэх, мэдээллийн санд хадгалах
үүрэгтэй. Хоёрдугаарт, хайлтын хэсэг нь хэрэглэгчийн асуултыг
боловсруулж хамааралтай хэсгүүдийг олж авах үүрэгтэй бөгөөд энэ хэсгийн
чанар нь системийн ерөнхий гүйцэтгэлд шийдвэрлэх нөлөө үзүүлдэг.
Гуравдугаарт, хариулт үүсгэх хэсэг нь олдсон хэсгүүдийг
контекст болгон ашиглан LLM-р хариулт гаргах үүрэгтэй.

\ \ \ \ \ \ \ \ \  \ \ \ \ RAG системийн хайлтын чанарыг дээшлүүлэх
чиглэлд сүүлийн жилүүдэд хэд хэдэн чухал дэвшил гарсан
\cite{chen2024benchmarking}. \textbf{Naive RAG} буюу энгийн RAG нь
хэрэглэгчийн асуултыг шууд векторт хөрвүүлж хайлт хийдэг анхны
арга бөгөөд хэрэглэхэд хялбар боловч нарийн асуултуудад чанар муудах
тал байдаг. Advanced RAG нь асуулт дахин боловсруулах, олон
дамжлагын хайлт, дахин эрэмбэлэх зэрэг нэмэлт алхмуудыг нэмж хайлтын
чанарыг дээшлүүлдэг. Modular RAG нь бүрэлдэхүүн хэсэг бүрийг
тусдаа модуль болгон боловсруулж, уян хатан байдлыг нэмэгдүүлсэн
хамгийн сүүлийн үеийн арга юм.

\ \ \ \ \ \ \ \ \  \ \ \ \ \ Hybrid хайлт нь утга санааны хайлт
болон түлхүүр үгийн хайлтыг хослуулсан арга бөгөөд RAG системийн
хайлтын чанарыг мэдэгдэхүйц дээшлүүлдэг болохыг судалгаанууд харуулсан
\cite{robertson2009bm25}. Утга санааны хайлт нь агуулгын хувьд ойрхон
мэдээллийг олоход давуу бол BM25 түлхүүр үгийн хайлт нь яг таарсан
нэр томьёо, тоо, хуулийн зүйл заалтын дугаар зэргийг олоход илүү үнэн
зөв байдаг. Эрх зүйн чиглэлд ``Иргэний хуулийн 450 дугаар зүйл'' гэх
мэт яг тодорхой лавлагаа хайхад BM25 арга нь утга санааны хайлтаас
илүү үр дүнтэй байдаг тул хоёуланг нь хослуулах нь зайлшгүй
шаардлагатай.

\ \ \ \ \ \ \ \ \  \ \ \ \ \Дахин эрэмбэлэх арга нь хайлтын
анхны үр дүнгүүдийг илүү нарийн загвараар дахин үнэлж хамгийн
хамааралтайг нь сонгодог \cite{nogueira2019passage}. Reciprocal Rank
Fusion (RRF) арга нь утга санааны болон түлхүүр үгийн хайлтын
үр дүнгүүдийг хялбар боловч үр дүнтэй аргаар нэгтгэдэг бөгөөд
тусгай сургалт шаардалгүйгээр сайн гүйцэтгэл үзүүлдэг давуу талтай
\cite{cormack2009reciprocal}.

\clearpage
\chapter{Технологийн судалгаа}
\markboth{}{}
\pagestyle{fancyplain}

\section{Системийн архитектурын ерөнхий бүтэц}

\ \ \ \ \ \ \ \ \  \ \ \ \ Энэхүү дипломын ажлын хүрээнд боловсруулж буй
хуулийн чатбот системийн технологийн сонгохдоо гүйцэтгэл,
хөгжүүлэлтийн хялбар байдал, Монгол хэлний дэмжлэг, ирээдүйд өргөтгөх
боломж болон зардлын үр ашиг гэсэн таван үндсэн шалгуурыг баримталсан.
Систем нь хэрэглэгчтэй харьцах вэб хэрэгсэл, хүсэлт боловсруулах серверийн
хэсэг, хиймэл оюун ухааны хайлт болон хариулт үүсгэх хэсэг, өгөгдөл
хадгалах хэсэг гэсэн дөрвөн давхаргаас бүрдэнэ \cite{richards2020fundamentals}.

\ \ \ \ \ \ \ \ \  \ \ \ \ Систем нь pnpm monorepo бүтэцтэй
хэрэгжсэн бөгөөд \texttt{apps/web} (Next.js вэб хэрэгсэл),
\texttt{apps/api} (Fastify серверийн хэсэг), \texttt{apps/worker}
(өгөгдөл боловсруулах ажилчин) болон \texttt{packages/shared}
(хуваалцсан төрөл, схем) гэсэн дөрвөн үндсэн модулиас бүрдэнэ.
Энэхүү бүтэц нь TypeScript төрлүүд болон баталгаажуулалтын схемийг
бүх модулиуд хооронд хуваалцах боломжийг олгодог тул алдааг
эрт илрүүлэх, кодын давталтыг багасгах ач холбогдолтой.

\section{Next.js — Вэб хэрэгслийн технологи}

\ \ \ \ \ \ \ \ \  \ \ \ \ \ Next.js нь Vercel компанийн хөгжүүлсэн
React суурилсан вэб хэрэгслийн фреймворк бөгөөд серверт болон
хэрэглэгч талд зэрэг ажиллах (Server-Side Rendering, SSR болон
Client-Side Rendering, CSR) боломжтой \cite{nextjs2024}. 2016 онд
анх гарсан энэхүү фреймворк өнөөдөр дэлхийн хамгийн өргөн
хэрэглэгддэг вэб фреймворкуудын нэг болж, 2024 оны байдлаар
npm-д долоо хоногт 6 сая гаруй удаа татагддаг болжээ.

\ \ \ \ \ \ \ \ \  \ \ \ \ Next.js 14 хувилбарт нэвтрүүлсэн
App Router архитектур нь хуудас бүрийн өгөгдөл татах
логикийг серверт шилжүүлж, анхны хуудас ачааллын хурдыг мэдэгдэхүйц
дээшлүүлдэг \cite{nextjs2024}. Server Components ашиглан чатботын
яриа хадгалах болон мессежийн жагсаалтыг серверт боловсруулж,
хэрэглэгч зөвхөн бэлэн HTML дамжуулах нь хэрэглэгчийн туршлагыг
сайжруулдаг. Мөн Server-Sent Events (SSE) буюу дамжуулалтын
горимыг дэмждэг тул LLM-ийн хариултыг токен токеноор дамжуулан
харуулах боломжтой — энэ нь хэрэглэгчид урт хариултыг хүлээх явцад
текст аажмаар гарч ирэхийг харах боломжийг олгодог.

\ \ \ \ \ \ \ \ \  \ \ \ \ \ React нь Next.js-ийн суурь болсон
хэрэглэгчийн интерфейсийн номын сан бөгөөд компонент суурилсан
архитектурт тулгуурладаг \cite{react2024}. Чатботын мессежийн жагсаалт,
оролтын хэсэг, эх сурвалжийн карт зэрэг UI элемент бүрийг тусдаа
компонент болгон хөгжүүлж, дахин ашиглах боломжтой болгосон.
\textbf{Tailwind CSS} ашиглан загварчлал хийснээр CSS файл тусдаа
бичих шаардлагагүй болж хөгжүүлэлтийн хурд нэмэгдсэн.

\begin{table}[h]
\centering
\caption{Frontend фреймворкуудын харьцуулалт}
\begin{tabular}{|l|c|c|c|c|}
\hline
\textbf{Шалгуур} & \textbf{Next.js} & \textbf{Vue/Nuxt} &
\textbf{Angular} & \textbf{SvelteKit} \\
\hline
SSR дэмжлэг        & Тийм   & Тийм   & Хэсэгчлэн & Тийм \\
SSE дэмжлэг        & Тийм   & Тийм   & Тийм       & Тийм \\
TypeScript         & Тийм   & Тийм   & Тийм       & Тийм \\
Экосистем          & Маш том & Том    & Том        & Жижиг \\
Сургах хугацаа     & Дунд   & Богино & Урт        & Богино \\
Гүйцэтгэл         & Өндөр  & Өндөр  & Дунд       & Маш өндөр \\
\hline
\end{tabular}
\label{tab:frontend_comparison}
\end{table}

\ \ \ \ \ \ \ \ \  \ \ \ \ Дээрх харьцуулалтаас харахад Next.js нь
SSR болон SSE-ийн нэгдмэл дэмжлэг, томоохон экосистем болон
TypeScript-тэй нягт нийцэх чадвараараа энэхүү системд хамгийн
тохиромжтой сонголт болсон. Vue/Nuxt хэдийгээр сургах хугацаа богино
боловч React-ийн экосистем болон нийгэмлэг (community) мэдэгдэхүйц
том байдаг тул тулгарах асуудлыг хурдан шийдвэрлэх боломж илүү
байдаг.

\section{Fastify — Серверийн хэсгийн технологи}

\ \ \ \ \ \ \ \ \  \ \ \ \ \textbf{Fastify} нь Node.js дээр суурилсан
өндөр гүйцэтгэлтэй HTTP фреймворк бөгөөд 2016 онд анх гарч
одоогоор npm-д долоо хоногт 3 сая гаруй удаа татагддаг болжээ
\cite{fastify2024}. Fastify-ийн үндсэн давуу тал нь
\textbf{JSON Schema}-д суурилсан хүсэлт болон хариултын
баталгаажуулалт бөгөөд энэ нь гаднаас ирэх буруу өгөгдлийг
серверийн логикт хүрэхээс өмнө шүүж зогсоодог. Манай системд
хэрэглэгчийн чат хүсэлтийн \texttt{message} талбарын урт, төрөл,
хоосон эсэхийг энэ аргаар шалгадаг.

\ \ \ \ \ \ \ \ \  \ \ \ \ Гүйцэтгэлийн хувьд Fastify нь Express.js-тэй
харьцуулахад секундэд ойролцоогоор хоёр дахин илүү хүсэлт
боловсруулдаг болохыг олон тооны бие даасан benchmark туршилт
харуулсан \cite{fastify2024}. RAG pipeline нь нэг хүсэлт
боловсруулахад Vector DB хайлт, LLM API дуудлага зэрэг олон
асинхрон үйлдэл дараалан хийдэг тул серверийн хурд нийт
хариу хугацаанд шууд нөлөөлдөг. Тиймээс хурдан фреймворк
сонгох нь стратегийн хувьд чухал шийдвэр байсан.

\ \ \ \ \ \ \ \ \  \ \ \ \ \textbf{Drizzle ORM} нь PostgreSQL-тэй
ажиллахад ашигласан TypeScript-ийн хөнгөн өгөгдлийн сангийн
зуучлагч юм \cite{drizzle2024}. Prisma болон TypeORM-тэй
харьцуулахад Drizzle нь SQL-д ойр синтакстай, цөөн хийсвэрлэлтэй
тул гүйцэтгэлийн хувьд давуу байдаг. TypeScript-ийн төрлийн
дэмжлэг бүрэн байдаг тул өгөгдлийн сангийн схем болон
TypeScript интерфейсийг нэг газраас удирдах боломжтой.

\begin{table}[h]
\centering
\caption{Node.js HTTP фреймворкуудын гүйцэтгэлийн харьцуулалт}
\begin{tabular}{|l|c|c|c|c|}
\hline
\textbf{Шалгуур} & \textbf{Fastify} & \textbf{Express} &
\textbf{Hono} & \textbf{Koa} \\
\hline
Req/сек (benchmark) & $\sim$77,000 & $\sim$35,000 &
$\sim$100,000 & $\sim$55,000 \\
JSON Schema         & Суурилагдсан & Гуравдагч & Хэсэгчлэн & Үгүй \\
TypeScript          & Тийм   & Тийм   & Тийм   & Тийм \\
Plugin систем       & Тийм   & Үгүй   & Үгүй   & Үгүй \\
Бичиг баримт        & Сайн   & Маш сайн & Сайн  & Дунд \\
Pino logging        & Суурилагдсан & Гуравдагч & Үгүй & Үгүй \\
\hline
\end{tabular}
\label{tab:backend_comparison}
\end{table}

\ \ \ \ \ \ \ \ \  \ \ \ \ Hono фреймворк Fastify-аас хурдан байдаг
хэдий ч plugin систем болон JSON Schema-ийн суурилагдсан дэмжлэг
байхгүй нь энэхүү системийн шаардлагад нийцэхгүй байсан. Express.js
нь хамгийн өргөн хэрэглэгддэг боловч гүйцэтгэл болон TypeScript-ийн
нэгдмэл дэмжлэгийн хувьд Fastify-д хоцордог.

\section{PostgreSQL болон pgvector}

\ \ \ \ \ \ \ \ \  \ \ \ \ \textbf{PostgreSQL} нь дэлхийн хамгийн
хөгжингүй нээлттэй эхийн харилцааны өгөгдлийн сан бөгөөд 1996 оноос
хойш тасралтгүй хөгжиж, өнөөдөр дэлхийн томоохон компаниудын дийлэнх
нь ашиглах болжээ \cite{postgresql2024}. ACID (Atomicity, Consistency,
Isolation, Durability) шинж чанарыг бүрэн хангадаг тул санхүүгийн
гүйлгээ, эрх зүйн баримт бичиг зэрэг нарийвчлал шаарддаг
өгөгдлийг хадгалахад онцгой тохиромжтой.

\ \ \ \ \ \ \ \ \  \ \ \ \ \textbf{pgvector} нь PostgreSQL-д вектор
хадгалах болон хайлт хийх чадварыг нэмдэг нээлттэй эхийн extension
юм \cite{pgvector2024}. 2021 онд анх гарсан энэхүү extension нь
\texttt{vector(n)} өгөгдлийн төрлийг нэмж, cosine distance
(\texttt{<=>}), Euclidean distance (\texttt{<->}), inner product
(\texttt{<\#>}) гэсэн гурван хайлтын аргыг дэмждэг. Манай системд
\texttt{vector(1536)} төрлийг ашиглан OpenAI-н
\texttt{text- embedding-3-small} загварын 1536 хэмжээст вектор
хадгалдаг.

\ \ \ \ \ \ \ \ \  \ \ \ \ pgvector-ийн хамгийн чухал онцлог бол нэг
өгөгдлийн сан дотор \textbf{вектор хайлт болон SQL хайлтыг
нэгтгэх} боломж юм. Тухайлбал дараах SQL асуулгаар вектор
болон бүтэн текстийн хайлтыг нэгэн зэрэг гүйцэтгэж болно.

\ \ \ \ \ \ \ \ \  \ \ \ \ pgvector нь хайлтын хурдыг дээшлүүлэхийн
тулд хоёр төрлийн индекс дэмждэг. \textbf{IVFFlat} индекс нь
өгөгдлийн санг тогтмол тооны бүлэгт хуваан хайлт хийдэг тул
санах ойн зардал бага боловч нарийвчлал бага зэрэг буурдаг.
\textbf{HNSW} (Hierarchical Navigable Small World) индекс нь
илүү нарийвчлалтай боловч санах ойн зардал өндөр байдаг
\cite{malkov2018hnsw}. Дипломын MVP шатанд 100,000 хүртэлх
баримтын хэмжээнд HNSW индекс нарийвчлалын хувьд илүү
тохиромжтой.

\begin{table}[h]
\centering
\caption{Вектор өгөгдлийн сангийн харьцуулалт}
\begin{tabular}{|l|c|c|c|c|}
\hline
\textbf{Шалгуур} & \textbf{pgvector} & \textbf{ChromaDB} &
\textbf{Pinecone} & \textbf{Weaviate} \\
\hline
SQL хайлт           & Тийм   & Үгүй   & Үгүй   & Үгүй \\
Нээлттэй эх         & Тийм   & Тийм   & Үгүй   & Тийм \\
Локал суулгалт      & Тийм   & Тийм   & Үгүй   & Тийм \\
ACID баталгаа       & Тийм   & Үгүй   & Үгүй   & Үгүй \\
Hybrid хайлт        & Тийм   & Үгүй   & Тийм   & Тийм \\
Тусдаа сервис       & Үгүй   & Тийм   & Тийм   & Тийм \\
Зардал              & Үнэгүй & Үнэгүй & Төлбөртэй & Үнэгүй \\
\hline
\end{tabular}
\label{tab:vectordb_comparison}
\end{table}

\section{ChromaDB — Вектор өгөгдлийн сан}

\ \ \ \ \ \ \ \ \  \ \ \ \ \textbf{ChromaDB} нь хиймэл оюун ухааны
програм хангамж хөгжүүлэхэд зориулсан нээлттэй эхийн вектор
өгөгдлийн сан бөгөөд 2022 онд Chroma Inc. компани анх гаргасан
\cite{chromadb2024}. Python болон JavaScript клиент номын сантай
бөгөөд локал файл систем дээр хялбар ажиллах боломжтой нь
хөгжүүлэлтийн орчинд туршилт хийхэд онцгой тохиромжтой болгодог.


\ \ \ \ \ \ \ \ \  \ \ \ \ Манай системд \texttt{VECTOR\_DB\_PROVIDER}
environment variable-р ChromaDB болон pgvector-ийн аль нэгийг
сонгох боломжтой болгосон. Хөгжүүлэлтийн орчинд ChromaDB,
production орчинд pgvector ашиглах нь зардал болон найдвартай
байдлын хувьд хамгийн оновчтой хослол болно.

\section{OpenAI API — Хэлний загвар болон Embedding}

\ \ \ \ \ \ \ \ \  \ \ \ \ \textbf{OpenAI}-н API нь манай системийн
хоёр үндсэн хиймэл оюун ухааны чадварыг хангадаг: текстийг
векторт хөрвүүлэх (embedding) болон хариулт үүсгэх (generation)
\cite{openai2024embedding}.

\ \ \ \ \ \ \ \ \  \ \ \ \ \textbf{text-embedding-3-small} загварыг
хуулийн текстийг векторт хөрвүүлэхэд ашигладаг. Энэхүү загвар нь
1536 хэмжээст вектор үүсгэдэг бөгөөд MTEB (Massive Text Embedding
Benchmark) олон хэлний шалгалтад өндөр үнэлгээ авсан
\cite{openai2024embedding}. Нэг сая токенд 0.02 ам.долларын
үнэтэй нь ижил чанарын загваруудтай харьцуулахад хамгийн хямд
сонголтуудын нэг юм. Монгол Кирилл үсгийг UTF-8 токенчлолоор
боловсруулдаг тул Монгол текстийн векторчлолд тохиромжтой.

\ \ \ \ \ \ \ \ \  \ \ \ \ \textbf{GPT-4o-mini} загварыг хэрэглэгчийн
асуултад хариулт үүсгэхэд ашигладаг \cite{openai2024gpt4o}.
128,000 токений контекст цонхтой энэхүү загвар нь олон хуулийн
баримт бичгийн хэсгийг нэгэн зэрэг контекст болгон хүлээн авах
чадвартай. GPT-4-тэй харьцуулахад зардал 10 дахин бага боловч
RAG контекстэд суурилсан хариулт гаргахад хангалттай чадвартай
болохыг туршилт харуулсан. Манай системийн хувьд зорилгот хариулт
нь загварын бие даасан мэдлэгт биш, харин retrieved chunk-уудад
тулгуурлан үүсгэгддэг тул загварын параметрийн тоо нь хариултын
чанарт шийдвэрлэх нөлөө үзүүлдэггүй.

\begin{table}[h]
\centering
\caption{LLM болон Embedding загваруудын харьцуулалт}
\begin{tabular}{|l|c|c|c|c|}
\hline
\textbf{Шалгуур} & \textbf{GPT-4o-mini} & \textbf{GPT-4o} &
\textbf{Claude 3.5} & \textbf{Gemini 1.5} \\
\hline
Контекст цонх    & 128K   & 128K   & 200K   & 1M \\
Монгол хэл       & Сайн   & Сайн   & Сайн   & Дунд \\
Зардал (1M tok)  & \$0.15 & \$5.00 & \$3.00 & \$0.35 \\
Хариу хугацаа    & Хурдан & Дунд   & Дунд   & Хурдан \\
RAG чанар        & Сайн   & Маш сайн & Маш сайн & Сайн \\
API хүртээмж     & Өндөр  & Өндөр  & Өндөр  & Дунд \\
\hline
\end{tabular}
\label{tab:llm_comparison}
\end{table}

\section{RAG архитектурын технологиуд}

\ \ \ \ \ \ \ \ \  \ \ \ \ RAG системийн гол бүрэлдэхүүнүүд болох
текстийг хэсэгчлэх, векторчлох, хайх болон дахин эрэмбэлэх
үйлдлүүдэд тусгай технологи болон номын санг ашигласан.

\ \ \ \ \ \ \ \ \  \ \ \ \ \textbf{Tiktoken} нь OpenAI-н токенчлолын
номын сан бөгөөд текстийг OpenAI загваруудын токенд хувааж
тоолох боломжийг олгодог \cite{tiktoken2024}. Манай системд
512 токений хэсэг үүсгэхдээ яг тодорхой токений тоог тоолохын
тулд tiktoken ашигладаг. Энэ нь тооцоолсон хэсгийн хэмжээ болон
бодит токений тоо таарахгүй байх эрсдэлийг арилгадаг.

\ \ \ \ \ \ \ \ \  \ \ \ \ \textbf{Cheerio} нь Node.js дахь HTML задлан
шинжлэгч бөгөөд jQuery-тэй ижил төстэй API-тай \cite{cheerio2024}.
\texttt{shuukh.mn} болон \texttt{legalinfo.mn}-с татсан HTML
баримтуудаас хуулийн үндсэн текст, гарчиг, огноо, хэргийн дугаар
зэргийг гаргахад ашигладаг. Cheerio нь \texttt{load()},
\texttt{\$('selector')}, \texttt{.text()} зэрэг энгийн командуудаар
HTML-ийн аль ч элементийн агуулгыг авах боломжтой тул хуулийн
баримт бичгийн бүтцэд тохирсон тусгай задлан шинжлэгч
хөгжүүлэхэд хялбар байсан.

\ \ \ \ \ \ \ \ \  \ \ \ \ \textbf{Zod} нь TypeScript-д зориулсан
схем тодорхойлолт болон баталгаажуулалтын номын сан юм
\cite{zod2024}. \texttt{packages/shared} модулд тодорхойлсон
Zod schema-г API-ийн оролт болон гаралт хоёуланд ашигладаг тул
frontend болон backend хооронд дамжих өгөгдлийн бүтэц үргэлж
тохирч байдаг. Тухайлбал \texttt{ChatRequestSchema} нь
\texttt{message} талбарыг 1-2000 тэмдэгтийн хязгаартай, мөрийн
төрлийн байхыг шаарддаг тул хэрэглэгчийн буруу оролтыг серверт
хүрэхээс өмнө шүүдэг.

\ \ \ \ \ \ \ \ \  \ \ \ \ \textbf{Node-cron} нь Node.js дахь cron
(тогтмол хуваарийн) ажлуудыг удирдах номын сан бөгөөд
\texttt{apps/worker} модулд шинэ хуулийн баримт бичгийг тогтмол
хуваарийн дагуу татах ажлыг удирдахад ашигладаг \cite{nodecron2024}.

\section{pnpm Monorepo болон Turborepo}

\ \ \ \ \ \ \ \ \  \ \ \ \ \pnpm нь npm болон yarn-д
өрсөлдөгч Node.js-ийн хамаарлын удирдлагын хэрэгсэл бөгөөд
тусгай ``content-addressable store'' аргаар давтагдах
package-уудыг нэг л удаа татаж хадгалдаг \cite{pnpm2024}.
npm-тэй харьцуулахад суулгах хурд мэдэгдэхүйц өндөр, зөвхөн
ашигласан package-уудад л хандах боломжтой байдаг ``strict
mode'' нь санамсаргүй хамаарлын асуудлаас хамгаалдаг.

\ \ \ \ \ \ \ \ \  \ \ \ \ \pnpm workspace нь monorepo
бүтцийн олон модулийг нэг репозиторид нэгтгэх боломжийг
олгодог \cite{pnpm2024}. \texttt{pnpm-workspace.yaml}
файлд тодорхойлсон бүх модуль хооронд package хуваалцах
боломжтой тул \texttt{packages/shared}-д тодорхойлсон
TypeScript төрлүүдийг \texttt{apps/api} болон \texttt{apps/web}
хоёулаа импортлон ашигладаг.

\ \ \ \ \ \ \ \ \  \ \ \ \ \Turborepo нь monorepo-д зориулсан
бүтээх болон туршилтын хэрэгсэл бөгөөд өөрчлөгдөөгүй модулийг
дахин бүтээхгүйгээр кэшлэн хадгалдаг \cite{turborepo2024}.
Жишээлбэл \texttt{packages/shared} өөрчлөгдөөгүй тохиолдолд
\texttt{apps/api}-г бүтээхэд shared-ийн бүтээлтийн үр дүнг
кэшнээс ашиглах тул бүтээлтийн хугацаа мэдэгдэхүйц богиносдог.

\begin{table}[h]
\centering
\caption{Package менежер болон Monorepo хэрэгслийн харьцуулалт}
\begin{tabular}{|l|c|c|c|c|}
\hline
\textbf{Шалгуур} & \textbf{pnpm+Turbo} & \textbf{npm+Lerna} &
\textbf{yarn+NX} & \textbf{Bun} \\
\hline
Суулгах хурд     & Маш хурдан & Удаан  & Хурдан & Маш хурдан \\
Кэш дэмжлэг     & Тийм       & Үгүй   & Тийм   & Хэсэгчлэн \\
TypeScript       & Тийм       & Тийм   & Тийм   & Тийм \\
Workspace        & Тийм       & Тийм   & Тийм   & Тийм \\
Тогтвортой байдал & Өндөр     & Дунд   & Өндөр  & Бага \\
Бичиг баримт     & Сайн       & Дунд   & Маш сайн & Хөгжиж байгаа \\
\hline
\end{tabular}
\label{tab:monorepo_comparison}
\end{table}

\section{Docker болон Container технологи}

\ \ \ \ \ \ \ \ \  \ \ \ \ \Docker нь програм хангамжийг
тусгаарлагдсан орчинд (container) ажиллуулах технологи бөгөөд
2013 онд анх гарч өнөөдөр дэлхийн хамгийн өргөн хэрэглэгддэг
container технологи болжээ \cite{docker2024}. Container нь
виртуал машинтай харьцуулахад илүү хөнгөн бөгөөд зөвхөн
програмын ажиллахад шаардлагатай хамаарлуудыг агуулдаг.

\ \ \ \ \ \ \ \ \  \ \ \ \ Манай системийн бүх бүрэлдэхүүнүүд
Docker Compose-р удирддаг тав тусдаа container дотор
ажилладаг. \texttt{web} (Next.js, port 3000),
\texttt{api} (Fastify, port 3001), \texttt{worker}
(Crawler/Ingestion), \texttt{postgres} (PostgreSQL + pgvector,
port 5432), \texttt{chroma} (ChromaDB, port 8000) гэсэн таван
сервис бие биеэсээ тусгаарлагдан ажилладаг тул нэг сервисийн
алдаа бусдад нөлөөлөхгүй.

\ \ \ \ \ \ \ \ \  \ \ \ \ \Docker volume} нь container
дахин эхлүүлэгдэх үед өгөгдөл алдагдахаас хамгаалдаг.
\texttt{pgdata} volume нь PostgreSQL-ийн бүх өгөгдлийг
container-ийн гадна хадгалдаг тул container устгагдсан ч
өгөгдөл хэвээр үлддэг. Мөн health check
аргаар postgres container бэлэн болсны дараа api
container эхлүүлэх дарааллыг баталгаажуулдаг.



\clearpage
\clearpage
\clearpage
\chapter{Системийн шаардлага}
\markboth{}{}

\section{Шаардлага тодорхойлох арга зүй}

\ \ \ \ \ \ \ \ \  \ \ \ \ Шаардлага тодорхойлоход дараах арга
зүйг ашигласан. Нэгдүгээрт, хэрэглэгчийн судалгаа —
хуулийн мэдээлэл хайдаг иргэд болон хуульч мэргэжилтнүүдтэй
ярилцлага хийж тулгамдсан асуудлуудыг тодорхойлсон.
Хоёрдугаарт, өрсөлдөгч системийн шинжилгээ —
\texttt{huuli.tech}, LexisNexis, LawGPT зэрэг системүүдийн
онцлог, давуу болон сул талыг судалсан. Гуравдугаарт,
техникийн шинжилгээ — RAG архитектурын онцлогоос
үүдэлтэй техникийн хязгаарлалтуудыг тодорхойлсон.

\ \ \ \ \ \ \ \ \  \ \ \ \ Энэхүү системийн архитектурын үндсэн
шийдвэрүүд нь шаардлагуудын хамрах хүрээг тодорхойлдог.
Систем нь hybrid RAG буюу хосолсон мэдээлэл хайлтад
суурилсан чатбот системийг хэрэгжүүлнэ. Уг архитектур нь
утга санааны хайлт болон түлхүүр үгийн хайлтыг зэрэгцүүлэн
гүйцэтгэж, дахин эрэмбэлэлтийн аргаар хамгийн хамааралтай
хэсгүүдийг сонгон LLM-д дамжуулдаг. Систем нь хэрэглэгчийн
бүртгэл болон нэвтрэлтийн аргагүй бөгөөд нэвтрэлгүй
хэрэглэгч шууд үйлчилгээг ашиглах боломжтой байна. Мэдлэгийн
санд хэрэглэгдэх өгөгдлийг \texttt{shuukh.mn} болон
\texttt{legalinfo.mn} эх сурвалжаас урьдчилан татаж авсны
үндсэн дээр цэвэрлэлт, хэсэгчлэлт, векторчлол болон хадгалалтын
бүрэн автоматжуулагдсан pipeline-р боловсруулж системд нэмнэ.

\section{Функциональ шаардлага}

\subsection{Хэрэглэгчтэй харилцах үйлдлүүд}

\begin{enumerate}

  \item \textbf{FR-01 — Нэвтрэлгүй хандалт.}
        Хэрэглэгч ямар нэгэн бүртгэл үүсгэх, нэвтрэх
        үйлдэл хийлгүйгээр системийг шууд ашиглах
        боломжтой байна. Хэрэглэгчийн хувийн мэдээлэл
        цуглуулах, хадгалах шаардлагагүй бөгөөд
        хандалтыг хязгаарлалт хэрэгжүүлэхгүй.
        \textit{Тэргүүлэх зэрэглэл: Заавал.}

  \item \textbf{FR-02 — Асуулт оруулах, баталгаажуулах.}
        Хэрэглэгч Монгол хэлээр хуулийн асуулт бичиж
        илгээх боломжтой байна. Оролт нь 1–2000
        тэмдэгтийн хязгаартай байх бөгөөд хоосон
        оролтыг хаана. Монгол Кирилл текстийг Unicode
        NFC нормчлолоор стандартчилах, prompt injection
        халдлагаас хамгаалах шаардлагатай. Оролтын
        тэмдэгтийн тоог бодит цагт харуулж 2000
        хязгаарт хүрэхэд дүрслэлийн анхааруулга гаргана.
        \textit{Тэргүүлэх зэрэглэл: Заавал.}

  \item \textbf{FR-03 — Эх сурвалжид тулгуурласан хариулт
        авах.}
        Хэрэглэгч асуулт илгээхэд систем холбогдох
        хуулийн мэдээллийг хайж олоод тэдгээрт
        тулгуурлан хариулт үүсгэнэ. Хариулт нь зөвхөн
        баримт бичгийн санд байгаа мэдээлэлд суурилах бөгөөд
        системийн баримт бичгийн санд байхгүй мэдээллийг
        зохиож хариулахгүй. Хариулт бүрт ашигласан
        хуулийн эх сурвалжийн нэр, холбоос болон
        холбогдох хэсгийн хураангуйг автоматаар гаргаж
        харуулна.
        \textit{Тэргүүлэх зэрэглэл: Заавал.}

  \item \textbf{FR-04 — Олон эргэлтийн яриа.}
        Хэрэглэгч нэг session-ийн хүрээнд дараалсан байдлаар олон асуулт асууж харилцах боломжтой байна. Систем нь өмнөх асуулт, хариултуудыг харгалзан үзэж, тухайн ярианы утга агуулгыг хадгалан контекстэд нийцсэн хариу өгнө. Иймд хэрэглэгчийн “Тэр торгуулийн хэмжээ хэд вэ?” гэх мэт дутуу тодорхойлолттой асуултыг өмнөх яриатай уялдуулан зөв ойлгож тайлбарлах чадвартай байна.
        \textit{Тэргүүлэх зэрэглэл: Заавал.}

  \item \textbf{FR-05 — Эх сурвалж шалгах боломж.}
        Хэрэглэгч хариулт бүрийн доор харуулагдах
        эх сурвалжийн холбоосыг ашиглан хариултын
        үндэслэлийг шалгах боломжтой байна. Эх
        сурвалжийн холбоосыг дарахад \texttt{legalinfo.mn}
        эсвэл \texttt{shuukh.mn}-ийн тухайн баримт
        бичигт шилжих боломжтой байна.
        \textit{Тэргүүлэх зэрэглэл: Заавал.}

  \item \textbf{FR-06 — Хариултын урсгал дамжуулалт.}
        LLM-ээс үүсгэж буй хариултыг Server-Sent Events (SSE) технологийг ашиглан хэсэгчлэн дамжуулж, хэрэглэгчид хэсэгчлэн харуулах боломжтой байна. Энэ нь урт хариултыг хүлээх явцад системийн хариу үйлдлийг шууд мэдрэх боломжийг бүрдүүлж, хэрэглэгчийн туршлагыг сайжруулна.
        \textit{Тэргүүлэх зэрэглэл: Хүсэлтэй.}

  \item \textbf{FR-07 — Эрх зүйн анхааруулга харуулах.}
        Систем хариулт бүрийн дор ``Энэхүү мэдээлэл
        нь мэргэжлийн хуулийн зөвлөгөө биш тул
        мэргэжлийн хуульчтай зөвлөлдөхийг зөвлөж
        байна'' гэсэн анхааруулгыг заавал харуулна.
        \textit{Тэргүүлэх зэрэглэл: Заавал.}

\end{enumerate}

\subsection{Системийн дотоод үйлдлүүд}

\begin{enumerate}

  \item \textbf{FR-08 — Утгад суурилсан хайлт.}
        Систем хэрэглэгчийн асуултыг 1536 хэмжээст
        векторт хөрвүүлж баримт бичгийн сангаас cosine
        similarity хэмжилтээр утгын хувьд
        хамааралтай хуулийн хэсгүүдийг хайж олно.
        Хайлтын анхны үр дүнгээс top-20 хамгийн
        хамааралтай хэсгийг авна.
        \textit{Тэргүүлэх зэрэглэл: Заавал.}

  \item \textbf{FR-09 — Түлхүүр үгийн хайлт.}
        Систем хуулийн зүйл заалтын дугаар, хэргийн
        дугаар болон яг таарсан нэр томьёогоор хайлт
        хийх боломжтой байна. PostgreSQL-ийн суурилагдсан
        full-text хайлтын боломжийг ашиглан BM25
        алгоритмаар үр дүнг эрэмбэлнэ.
        \textit{Тэргүүлэх зэрэглэл: Заавал.}

  \item \textbf{FR-10 — Хосолсон хайлт.}
        Утга санааны хайлт болон түлхүүр үгийн хайлтын
        үр дүнгүүдийг Reciprocal Rank Fusion алгоритмаар
        нэгтгэж давхардсан баримтуудыг цэвэрлэн
        top-20 нэгдсэн жагсаалт үүсгэнэ.
        \textit{Тэргүүлэх зэрэглэл: Заавал.}

  \item \textbf{FR-11 — Дахин эрэмбэлэлт.}
        Хосолсон хайлтаар олдсон top-20 хэсгүүдийг
        асуулт болон хэсэг бүрийн хамаарлыг нарийн
        үнэлэх cross-encoder загвараар дахин эрэмбэлж
        хамгийн хамааралтай top-5--8 хэсгийг сонгоно.
        Дахин эрэмбэлэлт нь зөвхөн vector зайн
        хэмжилтэд суурилсан хайлтаас илүү нарийвчлалтай
        эцсийн үр дүн гаргах боломжийг олгоно.
        \textit{Тэргүүлэх зэрэглэл: Заавал.}

  \item \textbf{FR-12 — Хариулт үүсгэх.}
        Дахин эрэмбэлэлтээр сонгогдсон top-5--8
        хэсгүүдийг системийн заавар болон асуултын
        түүхтэй нэгтгэн GPT-4o-mini загварт дамжуулж
        хариулт үүсгэнэ. Системийн заавар нь загварыг
        зөвхөн өгөгдсөн хэсгүүдэд тулгуурлан хариулахыг
        шаарддаг бөгөөд баримт бичгийн сангаас олдоогүй
        мэдээллийг зохиохыг хориглоно.
        \textit{Тэргүүлэх зэрэглэл: Заавал.}

  \item \textbf{FR-13 — Өгөгдөл боловсруулах pipeline.}
        Хөгжүүлэгч \texttt{shuukh.mn} болон
        \texttt{legalinfo.mn}-с гараар татаж авсан
        HTML баримтуудыг системд оруулахад дараах
        үйлдлүүд автоматаар гүйцэтгэгдэнэ: HTML
        цэвэрлэлт, Монгол текстийн Unicode NFC
        нормчлол, хуулийн зүйл заалтын хилийг таньж
        512 токений хэсэгт 64 токений давхцалтайгаар
        хуваах, \texttt{text- embedding-3-small}
        загвараар векторчлох, pgvector болон PostgreSQL-д
        хадгалах.
        \textit{Тэргүүлэх зэрэглэл: Заавал.}

  \item \textbf{FR-14 — Өгөгдлийн давхардал шалгах.}
        Системд өгөгдөл оруулахдаа URL-аар давхардлыг
        шалгаж аль хэдийн индексжсэн баримтыг дахин
        боловсруулахгүй байна. Давхардсан тохиолдолд
        upsert аргаар шинэ агуулгаар шинэчлэнэ.
        \textit{Тэргүүлэх зэрэглэл: Заавал.}

\end{enumerate}

\section{Функциональ бус шаардлага}


\subsection{Гүйцэтгэлийн шаардлага}

\begin{enumerate}

  \item \textbf{NFR-01 — Хариу хугацаа.}
        Chat endpoint-ийн p95 latency буюу хүсэлтүүдийн
        95 хувь нь 10 секундийн дотор хариулт авна.
        Хосолсон хайлт болон дахин эрэмбэлэлт нийт
        2 секундийн дотор дуусна. Дамжуулалтын горимд
        анхны токен 3 секундийн дотор гарч эхэлнэ.
        \textit{Хэмжих арга: Apache JMeter ашиглан
        ачааллын туршилт.}

  \item \textbf{NFR-02 — Зэрэгцээ хүсэлт боловсруулах.}
        Систем нэгэн зэрэг 10 хүртэлх хэрэглэгчийн
        хүсэлтийг гүйцэтгэлийн алдагдалгүйгээр
        боловсруулах чадвартай байна. MVP шатны
        зорилтот хэрэглэгчдийн тоо 100 хүртэлх байна.
        \textit{Тэргүүлэх зэрэглэл: Заавал.}

  \item \textbf{NFR-03 — Хайлтын чанар.}
        Туршилтын 50 гаруй асуулт дээр эх сурвалжийн
        лавлагааны нарийвчлал 80 хувиас дээш байна.
        Хариулт зөвхөн дахин эрэмбэлэлтээр сонгогдсон
        хэсгүүдэд тулгуурласан байх нь Faithfulness
        metric-ээр хэмжигдэнэ.
        \textit{Хэмжих арга: RAGAS framework.}

  \item \textbf{NFR-04 — Өгөгдөл боловсруулах хурд.}
        Өгөгдөл боловсруулах pipeline нь 100 баримт
        бичгийг 30 минутын дотор боловсруулж мэдлэгийн
        санд нэмэх чадвартай байна. Embedding үүсгэх
        үед OpenAI API-ийн rate limit-ийг харгалзан
        batch боловсруулалтыг ашиглана.
        \textit{Тэргүүлэх зэрэглэл: Хүсэлтэй.}

\end{enumerate}

\subsection{Аюулгүй байдлын шаардлага}

\begin{enumerate}

  \item \textbf{NFR-05 — Хүсэлтийн хязгаарлалт.}
        IP хаяг тутамд минутанд 10-аас илүү хүсэлт
        ирвэл 429 ``Too Many Requests'' хариулт буцаана.
        Хязгаарлалтыг давсан хэрэглэгчид retry хийх
        хугацааг хариулт дотор зааж өгнө.
        \textit{Тэргүүлэх зэрэглэл: Заавал.}

  \item \textbf{NFR-06 — Оролтын цэвэрлэгээ.}
        Хэрэглэгчийн бүх оролтыг Zod schema-р
        баталгаажуулж, prompt injection халдлагын
        нийтлэг хэв маягуудыг шүүнэ. API түлхүүр
        болон нууц мэдээллийг environment variable-д
        хадгалж код дотор байршуулахгүй байна.
        \textit{Тэргүүлэх зэрэглэл: Заавал.}

  \item \textbf{NFR-07 — Эх сурвалжид хандах.}
        Өгөгдөл гараар татахдаа \texttt{shuukh.mn}
        болон \texttt{legalinfo.mn} серверт хэт их
        дарамт учруулахгүйн тулд хүсэлт хоорондын
        зайг 2 секундээс доошгүй байлгана.
        \textit{Тэргүүлэх зэрэглэл: Заавал.}

\end{enumerate}

\subsection{Найдвартай байдлын шаардлага}

\begin{enumerate}

  \item \textbf{NFR-08 — Бүртгэл хяналт.}
        Бүх хүсэлт, хариулт, latency болон ашигласан
        токений тоог Pino structured logger-р JSON
        хэлбэрт бүртгэнэ. Өгөгдөл боловсруулах
        ажлуудын статусыг \texttt{crawl\_logs}
        хүснэгтэд хадгалж дараагийн ажлуудад ашиглана.
        \textit{Тэргүүлэх зэрэглэл: Заавал.}

  \item \textbf{NFR-09 — Алдааны боловсруулалт.}
        OpenAI API боломжгүй болсон тохиолдолд
        хэрэглэгчид ойлгомжтой алдааны мэдэгдэл
        харуулна. Vector DB холболт тасарсан
        тохиолдолд автоматаар дахин холбогдох
        оролдлого хийнэ. Өгөгдөл боловсруулах явцад
        гарсан алдааг \texttt{crawl\_logs} хүснэгтэд
        бүртгэж хадгална.
        \textit{Тэргүүлэх зэрэглэл: Заавал.}

  \item \textbf{NFR-10 — Өгөгдлийн тогтвортой байдал.}
        PostgreSQL болон pgvector-ийн өгөгдлийг Docker
        volume-р хадгалж container дахин эхлүүлэх
        үед алдагдахаас хамгаална.
        \textit{Тэргүүлэх зэрэглэл: Заавал.}

\end{enumerate}

\subsection{Хэрэглэхэд хялбар байдлын шаардлага}

\begin{enumerate}

  \item \textbf{NFR-11 — Responsive дизайн.}
        Вэб хэрэгсэл нь гар утас (320px-ээс), таблет
        болон компьютерийн дэлгэц дээр зохицон харагдана.
        Tailwind CSS-ийн breakpoint системийг ашиглан
        хэрэгжүүлнэ.
        \textit{Тэргүүлэх зэрэглэл: Хүсэлтэй.}

  \item \textbf{NFR-12 — Тэмдэгт тооны харуулалт.}
        Хэрэглэгч асуулт бичих явцад тэмдэгтийн тоог
        бодит цагт харуулна (жишээ: ``847 / 2000'').
        2000 хязгаарт хүрэхэд оролтын хайрцаг улаан
        болж анхааруулна.
        \textit{Тэргүүлэх зэрэглэл: Хүсэлтэй.}

  \item \textbf{NFR-13 — Ачааллын төлөв.}
        Систем хариулт боловсруулж байх явцад loading
        state харуулна. Дамжуулалтын горимд текст
        аажмаар гарч ирэх нь хэрэглэгчид системийн
        хариу үйлдлийг мэдрүүлдэг.
        \textit{Тэргүүлэх зэрэглэл: Хүсэлтэй.}

\end{enumerate}

\subsection{Өргөтгөх чадварын шаардлага}

\begin{enumerate}

  \item \textbf{NFR-14 — Vector DB солих боломж.}
        \texttt{VECTOR\_DB\_PROVIDER} environment
        variable-р ChromaDB болон pgvector-ийн аль
        нэгийг сонгох боломжтой байна. Retrieval
        сервис нь provider-ийн интерфейсийг хийсвэрлэн
        харьцдаг тул кодын өөрчлөлтгүйгээр солигдоно.
        \textit{Тэргүүлэх зэрэглэл: Хүсэлтэй.}

  \item \textbf{NFR-15 — LLM загвар солих боломж.}
        \texttt{OPENAI\_CHAT\_MODEL} болон
        \texttt{OPENAI\_EMBEDDING\_MODEL} environment
        variable-р загварыг тохиргооноос солих боломжтой
        байна. Generation сервис нь загварын API-тай
        хийсвэрлэгдсэн тул өөр провайдер руу шилжих
        боломжтой.
        \textit{Тэргүүлэх зэрэглэл: Хүсэлтэй.}

  \item \textbf{NFR-16 — Container-т суурилсан
        байршуулалт.}
        Систем бүхэлдээ Docker Compose-р ажиллах
        боломжтой байна. Орчны ялгаанаас үл хамааран
        ижил байдлаар ажиллах баталгааг хангана.
        \textit{Тэргүүлэх зэрэглэл: Заавал.}

\end{enumerate}

\section{Хязгаарлалт ба таамаглал}

\ \ \ \ \ \ \ \ \  \ \ \ \ Системийн хамрах хүрээг тодорхой
болгохын тулд дараах хязгаарлалт болон таамаглалуудыг
тодорхойлж баримтжуулсан. Эдгээр нь системийн шийдвэрлэхгүй
асуудлуудыг тодорхой болгодог бөгөөд ирээдүйн хөгжүүлэлтийн
чиглэлийг заана \cite{hull2011requirements}.

\begin{enumerate}

    \item \textbf{CON-01 — Өгөгдөл татах үйл ажиллагаа.}
        \texttt{shuukh.mn} болон \texttt{legalinfo.mn} цахим
        хуудаснаас өгөгдлийг эхний шатанд гараар татан авч
        бэлтгэнэ. Харин татаж авсан өгөгдлийг системд оруулах,
        боловсруулах pipeline нь автоматжуулсан байдлаар
        хэрэгжинэ.

  \item \textbf{CON-02 — Гуравдагч талын үйлчилгээний
        хамаарлал.}
        Систем нь OpenAI API-д шууд хамааралтай тул
        уг үйлчилгээ тасарсан тохиолдолд хариулт
        гаргах болон өгөгдөл боловсруулах боломжгүй
        болно. MVP шатанд нөөц LLM провайдер руу
        автоматаар шилжих арга хэрэгжүүлэхгүй.

  \item \textbf{CON-03 — Өгөгдлийн хамрах хүрээ.}
        Систем нь зөвхөн \texttt{shuukh.mn} болон
        \texttt{legalinfo.mn}-с гараар татаж авсан
        өгөгдөлд тулгуурлан хариулна. баримт бичгийн санд
        байхгүй хуулийн асуултад хангалттай хариулт
        өгч чадахгүй байж болзошгүй. 

  \item \textbf{CON-04 — Монгол хэлний боловсруулалтын
        хязгаарлалт.}
        Монгол хэлний тусгай tokenizer болон embedding
        загвар байхгүй тул OpenAI-н ерөнхий загварыг
        ашиглах шаардлагатай. Энэ нь Монгол хэлний
        нарийн утгын хайлтын чанарт тодорхой
        хязгаарлалт тавих магадлалтай болно.

  \item \textbf{CON-05 — Эрх зүйн таамаглал.}
        \texttt{shuukh.mn} болон \texttt{legalinfo.mn}
        нийтийн хандалттай хэвээр байна гэж тооцно.
        Гараар татаж авах үйлдэл нь техникийн болон
        ёс зүйн хувьд зөвшөөрөгдсөн гэж үзнэ.

  \item \textbf{CON-06 — Эрх зүйн зөвлөгөөний хязгаар.}
        Систем нь хуулийн мэдээлэл хайлтын хэрэгсэл
        бөгөөд мэргэжлийн хуулийн зөвлөгөө биш гэдгийг
        тодорхой заана. Системийн хариулт нь мэргэжлийн
        хуульчийн зөвлөгөөг орлохгүй гэдгийг
        хэрэглэгчид анхааруулна.

\end{enumerate}

\section{Use Case шинжилгээ}

\ \ \ \ \ \ \ \ \  \ \ \ \ Use case шинжилгээ нь системийн
хэрэглэгч болон системийн харилцан үйлдлийг тодорхойлоход
ашиглагдах бөгөөд Ivar Jacobson-ны аргачлалд тулгуурласан
\cite{jacobson1992object}. Манай системийн хувьд нэвтрэлгүй
хэрэглэгч гэсэн нэг төрлийн актор тодорхойлогдсон бөгөөд
дараах гурван үндсэн use case тодорхойлогдсон.

\begin{table}[h]
\centering
\caption{Use Case-уудын хураангуй}
\begin{tabular}{|l|l|l|}
\hline
\textbf{ID} & \textbf{Нэр} & \textbf{Тэргүүлэх зэрэглэл} \\
\hline
UC-01 & Хуулийн асуулт асуух   & Заавал \\
UC-02 & Яриа үргэлжлүүлэх     & Заавал \\
UC-03 & Эх сурвалж шалгах     & Заавал \\
\hline
\end{tabular}
\label{tab:usecase_summary}
\end{table}

\begin{figure}[h]
\centering
\includegraphics[width=0.8\textwidth]{usecasediagram.png}
\caption{Системийн Use Case диаграм}
\label{fig:usecase_diagram}
\end{figure}

Зураг \ref{fig:usecase_diagram}-т системийн үндсэн
актор болон use case-уудын хоорондын харилцааг
диаграм хэлбэрээр үзүүлэв.

\ \ \ \ \ \ \ \ \  \ \ \ \ \textbf{UC-01 — Хуулийн асуулт асуух}
нь системийн үндсэн use case юм. Нэвтрэлгүй хэрэглэгч chat
хэрэгсэлд Монгол хэлээр хуулийн асуулт бичиж илгээхэд
систем оролтыг баталгаажуулж, хосолсон хайлт гүйцэтгэж,
дахин эрэмбэлэлтээр хамгийн хамааралтай хэсгүүдийг сонгож,
LLM-р хариулт үүсгэж, эх сурвалжийн лавлагаа нэмж
хэрэглэгчид буцаана. Гол амжилтын нөхцөл бол хэрэглэгч
баримт бичгийн санд суурилсан хариулт болон баталгаажуулах
боломжтой лавлагаа авах явдал юм.

\ \ \ \ \ \ \ \ \  \ \ \ \ \textbf{UC-02 — Яриа үргэлжлүүлэх}
нь хэрэглэгч нэг session дотор өмнөх хариултад тулгуурлан
дахин асуулт асуух нөхцөлийг тодорхойлдог. Систем ярианы
өмнөх контекстийг LLM-д дамжуулснаар ``Тэр хуульд заасан
торгуулийн хэмжээ хэд вэ?'' гэх мэт дутуу тодорхойлолттой
асуултуудыг зөв тайлбарлах чадвартай байна.

\ \ \ \ \ \ \ \ \  \ \ \ \ \textbf{UC-03 — Эх сурвалж шалгах}
нь хэрэглэгч хариулт бүрийн доор харуулагдах source
card-уудыг ашиглан хариултын үндэслэлийг баталгаажуулах
нөхцөлийг тодорхойлдог. Хэрэглэгч эх сурвалжийн холбоосыг
дарж анхны баримт бичигт шилжих боломжтой байна.

\section{Шаардлагуудын хамаарлын матриц}

\ \ \ \ \ \ \ \ \  \ \ \ \ Шаардлагуудын хамаарлын матриц нь
функциональ шаардлага болон системийн бүрэлдэхүүнүүдийн
хоорондын хамаарлыг харуулж хэрэгжүүлэлтийн бүрэн
байдлыг шалгахад ашиглагддаг \cite{hull2011requirements}.

\begin{table}[h]
\centering
\caption{Шаардлага болон бүрэлдэхүүний хамаарлын матриц}
\begin{tabular}{|l|c|c|c|c|c|}
\hline
\textbf{Шаардлага} & \textbf{Web} & \textbf{API} &
\textbf{Worker} & \textbf{VectorDB} & \textbf{PostgreSQL} \\
\hline
FR-01 Нэвтрэлгүй хандалт  & \checkmark &            &            &            &            \\
FR-02 Асуулт оруулах      & \checkmark & \checkmark &            &            &            \\
FR-03 Хариулт авах        & \checkmark & \checkmark &            & \checkmark & \checkmark \\
FR-04 Олон эргэлт         & \checkmark & \checkmark &            &            & \checkmark \\
FR-05 Эх сурвалж шалгах   & \checkmark & \checkmark &            &            & \checkmark \\
FR-06 Дамжуулалт          & \checkmark & \checkmark &            &            &            \\
FR-07 Анхааруулга         & \checkmark &            &            &            &            \\
FR-08 Утга хайлт          &            & \checkmark &            & \checkmark &            \\
FR-09 Түлхүүр хайлт       &            & \checkmark &            &            & \checkmark \\
FR-10 Хосолсон хайлт      &            & \checkmark &            & \checkmark & \checkmark \\
FR-11 Дахин эрэмбэлэлт    &            & \checkmark &            &            &            \\
FR-12 Хариулт үүсгэх      &            & \checkmark &            &            &            \\
FR-13 Өгөгдөл боловсруулах &           &            & \checkmark & \checkmark & \checkmark \\
FR-14 Давхардал шалгах    &            &            & \checkmark & \checkmark & \checkmark \\
\hline
\end{tabular}
\label{tab:requirements_matrix}
\end{table}


\clearpage
\chapter{Хэрэгжүүлэлт}
\markboth{}{}

\section{Системийн ерөнхий бүтэц}

\ \ \ \ \ \ \ \ \  \ \ \ \ Энэхүү системийг pnpm workspace болон Turborepo дээр суурилсан monorepo байдлаар хэрэгжүүлсэн \cite{pnpm2024,turborepo2024}. Үндсэн хэрэгжүүлэлт нь \texttt{apps/web}, \texttt{apps/api}, \texttt{apps/worker}, \texttt{packages/shared} гэсэн дөрвөн хэсэгт төвлөрдөг. \texttt{packages/shared} модульд TypeScript төрөл, Zod schema, retrieval болон generation-т ашиглах тохиргооны тогтмолууд, мөн source URL болон текст боловсруулах туслах функцуудыг нэгтгэсэн. Ингэснээр frontend, backend, worker гурвын хооронд дамжих \texttt{ChatRequest}, \texttt{ChatResponse}, \texttt{Document}, \texttt{Chunk}, \texttt{Conversation} зэрэг өгөгдлийн бүтэц нэг мөр болж, кодын давхардал болон зөрүү багассан.

\ \ \ \ \ \ \ \ \  \ \ \ \ \texttt{apps/web} нь Next.js~14 дээр хэрэгжсэн хэрэглэгчийн интерфейс, \texttt{apps/api} нь Fastify-д суурилсан RAG сервер, \texttt{apps/worker} нь offline ingestion pipeline-г гүйцэтгэдэг тусдаа модуль юм. Репозиторид \texttt{packages/ui} гэсэн нэмэлт багц мөн байгаа боловч системийн үндсэн өгөгдлийн урсгал нь дээрх гурван application болон \texttt{packages/shared} хооронд төвлөрч байна. Ийм бүтэц сонгосноор crawler, retrieval, generation, chat UI, conversation persistence зэрэг үүргүүдийг салангид хөгжүүлэхийн зэрэгцээ нэг репозиторт нэгдсэнээрээ хамт тестлэх боломж бүрдсэн.

\begin{figure}[H]
\centering
\includegraphics[width=0.95\textwidth]
  {figures/monorepo_structure.png}
\caption{pnpm Monorepo-ийн модулиудын бүтэц болон хамаарал}
\label{fig:monorepo}
\end{figure}

\ \ \ \ \ \ \ \ \  \ \ \ \ Зураг \ref{fig:monorepo}-д модуль хоорондын хамаарлыг ерөнхийлөн үзүүлсэн. \texttt{packages/shared} нь API, Web, Worker гурвын дундах гэрээний үүргийг гүйцэтгэж, нэг модульд хийсэн schema эсвэл type-ийн өөрчлөлт бусад хэсэгт шууд тусах нөхцөлийг бүрдүүлдэг. Энэ нь ялангуяа chat request/response, document chunk metadata, conversation history зэрэг олон газар давтагддаг өгөгдлийн бүтэц дээр чухал ач холбогдолтой байв.

\section{Монгол хэлний текст боловсруулалт}

\ \ \ \ \ \ \ \ \  \ \ \ \ Системийн worker хэсэг нь \texttt{legalinfo} болон \texttt{shuukh} эх сурвалжаас татсан HTML хуудсуудыг parse хийхдээ эхлээд цэвэрлэгээ, дараа нь нормчилол, эцэст нь хэсэгчлэл гэсэн дарааллаар боловсруулдаг. Parser/cleaner шатанд HTML tag, илүүдэл whitespace, хоосон мөр, эвдэрхий control тэмдэгтүүдийг цэвэрлэж, Unicode NFC нормчилол хийснээр Монгол кирилл текстийг нэг стандарт төлөвт оруулдаг. Энэ алхам нь ижил харагдах боловч өөр code point-той тэмдэгтүүдээс үүдэх хайлтын алдааг бууруулахад чухал нөлөөтэй.

\ \ \ \ \ \ \ \ \  \ \ \ \ Хэсэгчлэлтийг \texttt{ChunkingService} хэрэгжүүлдэг бөгөөд эхлээд бүлэг, зүйл, заалт, дэд заалтын хилүүдийг илрүүлэх section-aware стратеги ашигладаг. Хэрэв баримт бичгийн бүтэц тодорхой байвал chunk бүрт \texttt{articleNo}, \texttt{clauseNo}, \texttt{subclauseNo}, \texttt{headingPath}, \texttt{sourceId}, \texttt{lawId}, \texttt{caseId}, \texttt{chunkType} зэрэг metadata хавсаргана. Харин бүтэц бүрэн тогтохгүй тохиолдолд sliding-window fallback ажиллаж, \texttt{CHUNK\_CONFIG}-т заасан 512 токенийн зорилтот хэмжээ, 64 токенийн давхцлыг ашиглан хөрш хэсгүүдийн логик холбоог хадгалдаг.

\begin{figure}[H]
\centering
\includegraphics[width=0.90\textwidth]
  {figures/chunking_strategy.png}
\caption{512 токенийн хэсэгчлэл болон 64 токенийн давхцлын стратеги}
\label{fig:chunking}
\end{figure}

\ \ \ \ \ \ \ \ \  \ \ \ \ Зураг \ref{fig:chunking}-д хэсэгчлэлийн зарчмыг ерөнхийлөн үзүүлсэн. Хуулийн нэг зүйл эсвэл заалтын төгсгөл дараагийн заалттай утгын холбоотой байх магадлал өндөр тул давхцал ашигласнаар retrieval шатанд чухал өгүүлбэр тасарч алдагдах эрсдэл буурдаг. Энэ нь Монгол хуулийн баримт бичгийн урт өгүүлбэр, шаталсан заалтын бүтэцтэй нийцсэн практик шийдэл болсон.

\section{Өгөгдлийн бааз ба вектор хадгалалт}

\ \ \ \ \ \ \ \ \  \ \ \ \ PostgreSQL нь системийн үндсэн өгөгдлийн сан бөгөөд чат болон баримт бичгийн мета өгөгдлийг хамтад нь хадгалдаг. Chat талд \texttt{conversations} ба \texttt{messages} хүснэгтүүд хэрэглэгч тус бүрийн харилцан яриаг бүртгэнэ. Assistant message-ийн metadata дотор \texttt{sources}, \texttt{relatedLaws}, \texttt{relatedCases}, \texttt{confidence}, \texttt{suggestedQuestions}, \texttt{usage} зэрэг бүтэцтэй өгөгдлийг хадгалдаг тул вэб интерфейс өмнөх хариултыг эх сурвалжтай нь хамт дахин ачаалах боломжтой болсон.

\ \ \ \ \ \ \ \ \  \ \ \ \ Ingestion талд \texttt{documents} ба \texttt{chunks} хүснэгтүүд татсан баримт бичиг болон тэдгээрийн хэсгүүдийг хадгална. Chunk бүрт текст, токены тоо, эх текст дэх \texttt{char\_start}/\texttt{char\_end} байрлал, \texttt{content\_hash}, embedding model, embedding dimensions, vector provider, мөн retrieval-д ашиглах metadata хадгалагддаг. Вектор хадгалалтын давхаргад provider abstraction ашигласан бөгөөд үндсэн provider нь \texttt{pgvector}, нэмэлт provider нь \texttt{Chroma} юм. Одоогийн тохиргоонд OpenAI-ийн \texttt{text-embedding-3-large} загварын 3072 хэмжээст embedding ашиглаж байгаа бөгөөд pgvector талд хэмжээ их үед \texttt{halfvec}-д суурилсан хайлтыг дэмжсэн. Ингэснээр ижил retrieval логикийг pgvector эсвэл Chroma дээр аль алинд нь ажиллуулах боломж бүрдсэн.

\section{Өгөгдөл боловсруулах Pipeline}

\ \ \ \ \ \ \ \ \  \ \ \ \ Offline ingestion pipeline-ийг \texttt{PipelineService} удирдаж, discovery $\rightarrow$ fetch $\rightarrow$ parse $\rightarrow$ chunk $\rightarrow$ embed $\rightarrow$ upsert гэсэн дарааллаар гүйцэтгэдэг. \texttt{LegalinfoCrawler} болон \texttt{ShuukhCrawler} нь эх сурвалж тус бүрийн discovery, fetch логикийг салангид хэрэгжүүлсэн тул pipeline нь ижил оркестрацитай ч source бүрт өөр crawler ашигладаг. Татсан баримт бичгээс \texttt{Document} бүтэц үүсгэж, дараа нь chunk үүсгэн embedding хийж, эцэст нь реляц өгөгдөл болон вектор хадгалалтад зэрэг шинэчилдэг.

\ \ \ \ \ \ \ \ \  \ \ \ \ Энэхүү pipeline-ийн практик давуу тал нь найдвартай ажиллагаанд оршино. Stage бүр дээр retry/backoff механизм, resume checkpoint, dead-letter JSONL бүртгэл, fail-fast босго, error-rate хяналт, мөн \texttt{content\_hash}-д суурилсан chunk deduplication ашигласан. Embedding service нь chunk-уудыг багцаар нь боловсруулж, OpenAI API ашиглах боломжгүй үед локал ONNX embedding рүү шилжих fallback замтай. Ийм шийдэл нь том хэмжээний crawl ажиллуулах үед тасалдлыг багасгаж, алдаатай баримтыг дараа нь тусад нь шинжлэх боломж олгосон.

\begin{figure}[H]
\centering
\includegraphics[width=0.95\textwidth]
  {figures/offline_ingestion_sequence.png}
\caption{Offline ingestion pipeline-ийн ерөнхий дараалал}
\label{fig:ingestion}
\end{figure}

\ \ \ \ \ \ \ \ \  \ \ \ \ Зураг \ref{fig:ingestion}-д worker pipeline-ийн дарааллыг ерөнхийлөн үзүүлсэн. Бодит хэрэгжүүлэлт дээр үүн дээр checkpoint хадгалалт, retry/backoff, dead-letter бүртгэл, duplicate chunk шүүлт зэрэг нэмэлт хамгаалалтын алхмууд ажиллаж байгаа нь зөвхөн дараалал гүйцэтгэхээс илүүтэйгээр найдвартай ingestion орчин бүрдүүлэхэд чиглэсэн.

\section{Hybrid RAG Pipeline}

\subsection{Хосолсон хайлт}

\ \ \ \ \ \ \ \ \  \ \ \ \ Хэрэглэгчийн асуулт орж ирэхэд chat route нь эхлээд retrieval plan байгуулж, асуултыг rewrite хийх, follow-up эсэхийг илрүүлэх, intent ангилах, шаардлагатай бол давуу эрхтэй хуулийн жагсаалт гаргах алхмуудыг гүйцэтгэдэг. Хэрэв client тал \texttt{conversationId} илгээсэн боловч message history явуулаагүй бол сервер өмнөх харилцан яриаг PostgreSQL-оос сэргээж retrieval-д ашигладаг. Ингэснээр товчилсон дараагийн асуултууд ч өмнөх контекст дээрээ тулгуурлан хайгдах боломжтой болсон.

\ \ \ \ \ \ \ \ \  \ \ \ \ Retrieval service нь асуултын embedding-ийг OpenAI эсвэл локал embedding provider-аар үүсгэж, дараа нь вектор хайлтыг pgvector эсвэл Chroma дээр ажиллуулна. Үүнтэй зэрэгцэн PostgreSQL full-text search дээр суурилсан BM25-тэй төстэй keyword хайлт явагдаж, хоёр жагсаалтыг Reciprocal Rank Fusion (RRF)-аар нэгтгэнэ. Үүний үр дүнд утгын ижил төстэй боловч яг ижил түлхүүр үггүй хэсгүүд, мөн нэр томьёо яг таарсан хэсгүүд хоёулаа нийлсэн candidate жагсаалт үүсдэг. Мөн retrieval шатанд эх сурвалжуудыг нэгтгэн \texttt{sources}, \texttt{relatedLaws}, \texttt{relatedCases} гэсэн бүтэцтэйгээр UI-д буцаах бэлтгэл хийгддэг.

\begin{figure}[H]
\centering
\includegraphics[width=0.90\textwidth]
  {figures/hybrid_search_rrf.png}
\caption{Хосолсон хайлт болон RRF нэгтгэлийн диаграм}
\label{fig:hybrid}
\end{figure}

\ \ \ \ \ \ \ \ \  \ \ \ \ Зураг \ref{fig:hybrid}-д hybrid retrieval-ийн үндсэн санааг үзүүлсэн. Практикт энэ нэгтгэл нь ганц vector search эсвэл ганц keyword search-оос шалтгаалах эрсдэлийг багасгаж, Монгол хуулийн нэр томьёо болон ерөнхий утгын хоёр төрлийн дохиог нэгэн зэрэг ашиглах боломж олгосон.

\subsection{Дахин эрэмбэлэлт ба хариулт үүсгэх}

\ \ \ \ \ \ \ \ \  \ \ \ \ Candidate chunk-уудыг нэгтгэсний дараа одоогийн хэрэгжүүлэлт \texttt{RerankerService}-ийн heuristic дахин эрэмбэлэлтийг ашигладаг. Энэ шатанд анхны vector score, query-term overlap, өндөр ач холбогдолтой signal term-ийн тааралт, анхны rank position, source төрөл зэрэг шалгууруудыг нэгтгэн шинэ score тооцоод эцсийн \texttt{top-N} context-ийг сонгодог. Ийм шийдэл нь cross-encoder ашиглаагүй боловч retrieval-ийн дараах шүүлтийн чанарыг мэдэгдэхүйц сайжруулж, илүү хамааралтай хэсгүүдийг generation шатанд дамжуулах нөхцөл бүрдүүлсэн.

\begin{figure}[H]
\centering
\includegraphics[width=0.90\textwidth]
  {figures/reranking_diagram.png}
\caption{Дахин эрэмбэлэлтийн үе шатны ерөнхий тайлбар}
\label{fig:rerank}
\end{figure}

\ \ \ \ \ \ \ \ \  \ \ \ \ Зураг \ref{fig:rerank} нь дахин эрэмбэлэлтийн санааг ерөнхийлөн үзүүлж байгаа бөгөөд бодит кодын хэрэгжүүлэлтэд дээр дурдсан heuristic reranker ажиллаж байна. Дахин эрэмбэлэгдсэн chunk-уудыг generation service хүлээн авч, ярианы хамааралтай түүх, шүүгдсэн related case мэдээлэл, системийн prompt-той нэгтгэн OpenAI chat completion дуудна. Grounded answer үүсгэхдээ \texttt{legalinfo} эх сурвалжийн chunk-уудыг давуу ашиглаж, шаардлагатай үед \texttt{shuukh} кейсүүдийг нэмэлт тайлбар эсвэл fallback зөвлөмж хэлбэрээр тусгадаг. Хэрэв контекст сул байвал LLM-д шууд найдахгүйгээр \texttt{no-info} эсвэл \texttt{fallback-general} горимын хариулт буцаах зам ажиллана.

\begin{figure}[H]
\centering
\includegraphics[width=0.95\textwidth]
  {figures/rag_pipeline_sequence.png}
\caption{Hybrid RAG pipeline-ийн ерөнхий sequence диаграм}
\label{fig:rag_sequence}
\end{figure}

\ \ \ \ \ \ \ \ \  \ \ \ \ Зураг \ref{fig:rag_sequence}-д асуултаас хариулт хүртэлх ерөнхий урсгалыг харуулсан. Одоогийн хэрэгжүүлэлт дээр generation-ийн үндсэн загвар нь \texttt{gpt-4o-mini} бөгөөд API нь Web UI руу нэг бүрэн JSON хариулт буцаадаг. Гэсэн ч зурагт үзүүлсэн шиг retrieval, rerank, generation, persistence гэсэн үндсэн шатлал нь бодит урсгалтай нийцэж байгаа.

\section{Fastify API хэрэгжүүлэлт}

\ \ \ \ \ \ \ \ \  \ \ \ \ Серверийн хэсэг нь Fastify framework дээр plugin хэлбэрийн зохион байгуулалттай хэрэгжсэн \cite{fastify2024}. Authentication, CORS, rate limiting, structured logging зэрэг middleware шинжтэй боломжуудыг тусдаа plugin хэлбэрээр салгасан нь route handler-уудыг төвөггүй, тестлэхэд хялбар болгосон. API-ийн shared contract нь \texttt{packages/shared}-ийн Zod schema дээр суурилдаг тул request validation болон runtime type safety-ийг нэг эх сурвалжаас хангаж байна.

\ \ \ \ \ \ \ \ \  \ \ \ \ \texttt{POST /v1/chat} endpoint нь системийн гол урсгалыг хэрэгжүүлдэг. Route handler нь эхлээд хэрэглэгчийн нэвтрэлтийг шалгаж, дараа нь request body-г \texttt{chatRequestSchema}-аар баталгаажуулна. \texttt{conversationId} ирсэн бол тухайн conversation тухайн хэрэглэгчид хамаарах эсэхийг шалгаж, шаардлагатай бол message history-г сэргээнэ. Үүний дараа retrieval plan байгуулж, search $\rightarrow$ rerank $\rightarrow$ generate дарааллыг ажиллуулаад conversation болон user/assistant message-үүдийг PostgreSQL-д хадгална. Default rate limit нь 60 секундэд 10 хүсэлт байхаар тохируулагдсан бөгөөд validation, authentication, ownership алдаанд тохирсон HTTP статусууд буцаадаг.

\begin{figure}[H]
\centering
\includegraphics[width=0.90\textwidth]
  {figures/post_chat_sequence.png}
\caption{POST /v1/chat endpoint-ийн ерөнхий урсгал}
\label{fig:chat_seq}
\end{figure}

\ \ \ \ \ \ \ \ \  \ \ \ \ Route-ийн JSON хариултад \texttt{conversationId}, \texttt{answer}, \texttt{sources}, \texttt{relatedLaws}, \texttt{relatedCases}, \texttt{confidence}, \texttt{sourcesUsed}, \texttt{suggestedQuestions}, \texttt{usage.latencyMs} зэрэг талбарууд ордог. Ийм бүтэц нь вэб интерфейст зөвхөн текстэн хариулт төдийгүй холбоотой хууль, кейс, эх сурвалж, тайлбарлах follow-up асуултуудыг зэрэг харуулах боломж олгодог. Зураг \ref{fig:chat_seq}-д энэхүү endpoint-ийн гол алхмуудыг нэгтгэн үзүүлсэн.

\section{Next.js вэб хэрэгжүүлэлт}

\ \ \ \ \ \ \ \ \  \ \ \ \ Вэб хэсэг нь Next.js~14-ийн App Router архитектур дээр хэрэгжсэн \cite{nextjs2024}. Chat UI-ийн үндсэн зангилаа нь \texttt{ChatShell} бөгөөд энэ компонент \texttt{ChatInput}, \texttt{MessageList}, \texttt{SourcesAccordion} зэрэг дэд компонентуудыг нэгтгэн ажиллуулдаг. Хэрэглэгчийн authentication төлөв, conversation list, selected conversation, message history, loading/error төлөвүүдийг \texttt{useChat} hook төвлөрүүлэн удирдаж байна.

\ \ \ \ \ \ \ \ \  \ \ \ \ \texttt{useChat} hook нь message илгээхэд эхлээд хэрэглэгчийн бичсэн текстийг UI дээр түр нэмж харуулж, дараа нь \texttt{sendChatMessage()} дуудлагаар API руу хүсэлт явуулдаг. Хариулт амжилттай ирвэл assistant message-ийг sources, related laws/cases, confidence, latency зэрэг дагалдах өгөгдөлтэй нь state-д нэмнэ. Үүнээс гадна \texttt{fetchConversations()} болон \texttt{fetchConversation()} функцууд нь өмнөх яриануудыг ачаалж, \texttt{retry()} нь алдаа гарсан хамгийн сүүлийн асуултыг дахин илгээдэг. Ийм төвлөрсөн state management хийснээр чат UI-ийн зан төлөв илүү энгийн, урьдчилан таамаглах боломжтой болсон.

\begin{figure}[H]
\centering
\includegraphics[width=0.90\textwidth]
  {figures/streaming_sse_sequence.png}
\caption{Web UI ба API хоорондын чат харилцааны ерөнхий sequence диаграм}
\label{fig:sse}
\end{figure}

\ \ \ \ \ \ \ \ \  \ \ \ \ Одоогийн кодын хэрэгжүүлэлт нь token streaming бус, нэг бүрэн JSON хариулт хүлээн аваад UI state-аа шинэчилдэг request-response загвартай. Тиймээс assistant message нь хариулт бүрэн ирсний дараа дэлгэцэнд нэмэгдэж, түүнтэй хамт эх сурвалж болон related мэдээллүүд гарч ирдэг. Зураг \ref{fig:sse} нь чат харилцааны sequence-ийг ерөнхийлөн харуулж байгаа бөгөөд бодит хувилбарт түүний streaming хэсгийг бүрэн JSON response-оор орлуулсан гэж үзэж болно.

\section{Docker байршуулалт}

\ \ \ \ \ \ \ \ \  \ \ \ \ Системийн container орчныг Docker Compose-оор тодорхойлсон \cite{docker2024}. Локал хөгжүүлэлтийн \texttt{docker-compose.yml} файлд \texttt{postgres} (pgvector extension-тэй), \texttt{chroma}, \texttt{api}, \texttt{worker}, \texttt{pgadmin} сервисүүд багтдаг. Эдгээр нь нэг bridge network дээр ажиллаж, \texttt{pgdata}, \texttt{chromadata}, \texttt{workerdata}, \texttt{pgadmin-data} volume-уудаар дамжин өгөгдлөө container-оос тусад нь хадгалдаг. \texttt{api} болон \texttt{worker} сервисүүдийн орчинд \texttt{VECTOR\_DB\_PROVIDER=pgvector} нь анхдагч утга боловч Chroma сервис зэрэг ажиллаж байгаа нь fallback эсвэл өөр provider турших боломжийг хадгалдаг.

\ \ \ \ \ \ \ \ \  \ \ \ \ Production override болох \texttt{docker-compose.prod.yml} файлд \texttt{web} сервис нэмэгдэж, \texttt{NEXT\_PUBLIC\_API\_URL} нь дотоод \texttt{api} контейнер руу чиглэнэ. Харин локал compose дээр \texttt{web} сервис заавал асахгүй, frontend-ийг тусад нь хөгжүүлэлтийн горимоор ажиллуулах боломж нээлттэй үлдсэн. \texttt{worker} контейнер нь scheduler горимоор pipeline-аа ажиллуулж, \texttt{postgres}, \texttt{api}, \texttt{chroma} сервисийн бэлэн байдлаас хамаарч асах байдлаар зохион байгуулагдсан.

\begin{figure}[H]
\centering
\includegraphics[width=0.85\textwidth]
  {figures/docker_deployment.png}
\caption{Docker Compose container-уудын бүтэц болон хамаарал}
\label{fig:docker}
\end{figure}

\ \ \ \ \ \ \ \ \  \ \ \ \ Зураг \ref{fig:docker}-д container хоорондын хамаарлыг ерөнхийлөн үзүүлсэн. Бодит deployment дээр локал болон production compose-ийн бүтэц бага зэрэг ялгаатай боловч API, Worker, PostgreSQL, Chroma гэсэн үндсэн сервисийн зааг өөрчлөгдөөгүй. Ийм байршуулалт нь development үед ingestion болон retrieval давхаргыг тусгаарлан турших, production үед бүрэн стекээр нь ажиллуулах хоёр хувилбарыг хоёуланг нь дэмждэг.
\chapter{Туршилт ба үр дүн}
\markboth{}{}

\section{Туршилтын зорилго ба хамрах хүрээ}

\ \ \ \ \ \ \ \ \  \ \ \ \ Энэхүү бүлэгт хэрэгжүүлсэн системийн үндсэн
шаардлагууд бодит ажиллагаанд хэрхэн хангагдаж байгааг шалгасан
туршилтын арга, үр дүнг нэгтгэн үзүүлэв. Туршилтын зорилго нь
нэгдүгээрт өгөгдөл цуглуулах болон индексжүүлэх pipeline тогтвортой
ажиллаж байгаа эсэх, хоёрдугаарт RAG pipeline хэрэглэгчийн асуултад
эх сурвалжид тулгуурласан хариулт үүсгэж байгаа эсэх, гуравдугаарт
API болон үндсэн сервисүүдийн автомат тестүүд амжилттай ажиллаж байгаа
эсэхийг баталгаажуулахад чиглэсэн.

\ \ \ \ \ \ \ \ \  \ \ \ \ Туршилтыг дөрвөн түвшинд гүйцэтгэсэн. Үүнд:
өгөгдөл боловсруулалтын туршилт, backend сервисийн нэгж тест, API
интеграцийн туршилт, хэрэглэгчийн асуултаар хийсэн гар ажиллагааны
чанарын шалгалт багтана. Туршилтын орчин нь локал хөгжүүлэлтийн
Docker Compose орчин бөгөөд PostgreSQL, ChromaDB, Fastify API болон
worker сервисүүдийг тус тусад нь ажиллуулж шалгасан.

\begin{table}[H]
\renewcommand{\tablename}{Хүснэгт}
\centering
\caption{Туршилтын хамрах хүрээ}
\begin{tabular}{|l|l|l|}
\hline
\textbf{Туршилтын төрөл} & \textbf{Шалгасан бүрэлдэхүүн} & \textbf{Зорилго} \\
\hline
Өгөгдөл боловсруулалт & Worker, crawler, chunker, embedding & Баримтыг индексжүүлэх \\
Нэгж тест & API service, query rewrite, generation & Логикийн зөв байдлыг шалгах \\
Интеграцийн тест & Chat route, conversation flow & API урсгалын ажиллагааг шалгах \\
Гар ажиллагааны тест & Web UI, RAG answer, sources & Хариултын чанарыг үнэлэх \\
\hline
\end{tabular}
\label{tab:test_scope}
\end{table}

\section{Өгөгдөл боловсруулалтын туршилт}

\ \ \ \ \ \ \ \ \  \ \ \ \ Offline ingestion pipeline-ийн ажиллагааг
\texttt{data/raw} хавтас дахь legalinfo болон shuukh эх сурвалжийн
HTML файлууд дээр туршсан. Туршилтын үед нийт 113 HTML файл бэлэн
байснаас 100 баримт бичгийг pipeline-р боловсруулж, текст цэвэрлэгээ,
chunk үүсгэлт, embedding тооцоолол болон ChromaDB-д upsert хийх
дарааллыг гүйцэтгэсэн. Үр дүнд нь 100 chunk үүсэж, бүгд 1536 хэмжээст
embedding-тэйгээр вектор санд амжилттай хадгалагдсан.

\begin{table}[H]
\renewcommand{\tablename}{Хүснэгт}
\centering
\caption{Өгөгдөл боловсруулалтын туршилтын үр дүн}
\begin{tabular}{|l|c|}
\hline
\textbf{Үзүүлэлт} & \textbf{Утга} \\
\hline
Бэлэн HTML файл & 113 \\
Боловсруулсан баримт & 100 \\
Үүссэн chunk & 100 \\
Embedding хийгдсэн chunk & 100 \\
Vector DB-д хадгалсан chunk & 100 \\
Алдаа & 0 \\
Амжилтын хувь & 100\% \\
Нийт хугацаа & 18.353 секунд \\
Дундаж хугацаа / баримт & 183.5 ms \\
\hline
\end{tabular}
\label{tab:ingestion_result}
\end{table}

\ \ \ \ \ \ \ \ \  \ \ \ \ Энэ үр дүнгээс харахад локал өгөгдөл дээр
pipeline-ийн үндсэн дараалал тасралтгүй ажиллаж, баримт бичгийг
вектор хайлтад бэлэн төлөвт оруулж чадсан. Гэсэн хэдий ч туршилтын
үед нэг баримтаас дунджаар нэг chunk үүссэн нь том хэмжээтэй хууль,
шүүхийн шийдвэрийн өгөгдөл дээр илүү нарийн section-aware chunking
тохиргоог дахин шалгах шаардлагатайг харуулсан. Өөрөөр хэлбэл ingestion
амжилттай болсон боловч chunk-ийн тоо болон хэмжээ нь хариултын чанарт
шууд нөлөөлөх тул цаашид бүрэн хэмжээний өгөгдөл дээр дахин тохируулга
хийх хэрэгтэй.

\section{Автомат тестийн үр дүн}

\ \ \ \ \ \ \ \ \  \ \ \ \ Backend API талын автомат тестийг Vitest
ашиглан ажиллуулсан. API тестүүд query rewrite, scope classifier,
chat workflow, generation service зэрэг RAG pipeline-ийн гол логикийг
хамарсан. Туршилтын үр дүнд 6 test file-оос 4 нь амжилттай ажиллаж,
нийт 17 test давсан бөгөөд 9 test нь todo/skipped төлөвтэй байсан.

\begin{table}[H]
\renewcommand{\tablename}{Хүснэгт}
\centering
\caption{API автомат тестийн үр дүн}
\begin{tabular}{|l|c|c|c|}
\hline
\textbf{Тестийн файл} & \textbf{Pass} & \textbf{Skipped/Todo} & \textbf{Төлөв} \\
\hline
query-rewrite.service.test.ts & 6 & 0 & Амжилттай \\
scope-classifier.service.test.ts & 1 & 0 & Амжилттай \\
chat-workflow.service.test.ts & 3 & 0 & Амжилттай \\
generation.service.test.ts & 7 & 0 & Амжилттай \\
chat.route.test.ts & 0 & 5 & Хойшлуулсан \\
retrieval.service.test.ts & 0 & 4 & Хойшлуулсан \\
\hline
\textbf{Нийт} & \textbf{17} & \textbf{9} & \textbf{Pass} \\
\hline
\end{tabular}
\label{tab:api_test_result}
\end{table}

\ \ \ \ \ \ \ \ \  \ \ \ \ Generation service-ийн тестүүд нь систем
хангалттай эх сурвалжгүй үед хийсвэр хариулт зохиохгүй байх, хөдөлмөрийн
маргаан болон зам тээврийн нөхцөлд fallback зөвлөгөө өгөх, follow-up
асуултад өмнөх түүхийг ашиглах зэрэг тохиолдлуудыг шалгасан. Энэ нь
FR-02, FR-03 болон FR-11 шаардлагуудын үндсэн логик зөв ажиллаж байгааг
харуулж байна.

\ \ \ \ \ \ \ \ \  \ \ \ \ Worker талын тестийг мөн Vitest-р ажиллуулахад
19 test-ээс 8 нь амжилттай, 5 нь skipped/todo, 6 нь failed төлөвтэй
гарсан. Алдааны шалтгааныг задлан үзвэл ingest job-ийн 3 test нь
локал PostgreSQL-ийн \texttt{5433} порт дээр холболт үүсээгүйгээс
уналсан. Харин chunker-ийн 3 test нь шинэчлэгдсэн chunking логик
болон тестийн хүлээгдэж буй token/metadata утгуудын зөрүүнээс үүдсэн.

\begin{table}[H]
\renewcommand{\tablename}{Хүснэгт}
\centering
\caption{Worker автомат тестийн үр дүн}
\begin{tabular}{|l|c|c|c|l|}
\hline
\textbf{Тестийн файл} & \textbf{Pass} & \textbf{Fail} & \textbf{Skipped} & \textbf{Тайлбар} \\
\hline
normalizer.test.ts & 5 & 0 & 0 & Амжилттай \\
chunker.test.ts & 3 & 3 & 0 & Chunk metadata зөрсөн \\
ingest.test.ts & 0 & 3 & 0 & PostgreSQL холболтгүй \\
shuukh.crawler.test.ts & 0 & 0 & 5 & Хойшлуулсан \\
\hline
\textbf{Нийт} & \textbf{8} & \textbf{6} & \textbf{5} & \textbf{Хэсэгчлэн} \\
\hline
\end{tabular}
\label{tab:worker_test_result}
\end{table}

\ \ \ \ \ \ \ \ \  \ \ \ \ Иймээс автомат тестийн түвшинд API-ийн
үндсэн logic тогтвортой байгаа боловч worker-ийн chunker болон ingest
test-үүдийг одоогийн implementation-тэй нийцүүлэн шинэчлэх шаардлагатай
гэж дүгнэв. Энэ нь системийн бодит ingestion ажилласан үр дүнг үгүйсгэхгүй,
харин regression test-ийн хүлээлт шинэ хэрэгжүүлэлттэй бүрэн таараагүй
байгааг харуулж байна.

\section{Гүйцэтгэлийн туршилт}

\ \ \ \ \ \ \ \ \  \ \ \ \ Системийн гүйцэтгэлийг offline pipeline
болон online query pipeline гэсэн хоёр хэсэгт үнэлсэн. Offline талд
100 баримтыг 18.353 секундэд боловсруулсан нь нэг баримтад дунджаар
183.5 ms зарцуулсан гэсэн үг юм. Online query pipeline-ийн хувьд
сервис бүрийн хэмжилтээс vector search 50--100 ms, BM25 search
20--50 ms, RRF fusion 5--10 ms, reranking 10--30 ms, LLM generation
1--3 секунд орчим хугацаа зарцуулж байсан.

\begin{table}[H]
\renewcommand{\tablename}{Хүснэгт}
\centering
\caption{Query pipeline-ийн latency хэмжилт}
\begin{tabular}{|l|c|}
\hline
\textbf{Pipeline үе шат} & \textbf{Latency} \\
\hline
Vector search & 50--100 ms \\
BM25 keyword search & 20--50 ms \\
RRF fusion & 5--10 ms \\
Reranking & 10--30 ms \\
LLM generation & 1--3 секунд \\
Нийт query хугацаа & 1.2--3.5 секунд \\
\hline
\end{tabular}
\label{tab:latency_result}
\end{table}

\ \ \ \ \ \ \ \ \  \ \ \ \ Дээрх хэмжилтээс харахад retrieval болон
reranking шатны хугацаа нийт хариу боловсруулах хугацааны бага хэсгийг
эзэлж, LLM generation хамгийн их хугацаа зарцуулж байна. Иймээс
хэрэглэгчийн мэдрэгдэх хурдыг сайжруулах хамгийн үр дүнтэй арга нь
token streaming, response cache болон prompt/context-ийн хэмжээг
оновчлох явдал гэж үзсэн. NFR-02-т заасан 30 секундээс бага хугацаанд
бүрэн хариу өгөх шаардлагыг хэмжилтийн орчинд хангаж байсан бөгөөд
NFR-03-т заасан эхний хариу ирэх хугацааг сайжруулахад SSE streaming
бүрэн идэвхжүүлэх шаардлагатай.

\section{Функциональ шалгалтын үр дүн}

\ \ \ \ \ \ \ \ \  \ \ \ \ Гар ажиллагааны функциональ шалгалтад
хуулийн зүйл заалт асуух, ерөнхий эрх зүйн нөхцөл тайлбарлуулах,
өмнөх асуултаа үргэлжлүүлэх, эх сурвалж шалгах, мэдээлэл олдоогүй
үед fallback хариулт буцаах зэрэг хэрэглээний тохиолдлуудыг хамруулсан.
Эдгээр тохиолдлууд нь шаардлагын бүлэгт тодорхойлсон FR-01--FR-06
болон retrieval-тэй холбоотой FR-07--FR-11 шаардлагуудтай шууд
холбоотой.

\begin{table}[H]
\renewcommand{\tablename}{Хүснэгт}
\centering
\caption{Функциональ шаардлагын шалгалтын хураангуй}
\begin{tabular}{|l|l|l|}
\hline
\textbf{Шаардлага} & \textbf{Шалгасан үйлдэл} & \textbf{Үр дүн} \\
\hline
FR-01 & Асуултын оролт, validation & Хангасан \\
FR-02 & Эх сурвалжид суурилсан хариулт & Хангасан \\
FR-03 & Олон эргэлтийн яриа & Хангасан \\
FR-04 & Source URL харуулах & Хангасан \\
FR-05 & Streaming хариулт & Хэсэгчлэн \\
FR-06 & Эрх зүйн анхааруулга & Хангасан \\
FR-07--FR-10 & Vector, keyword, RRF, rerank & Хангасан \\
FR-11 & LLM хариулт үүсгэх & Хангасан \\
FR-12 & Давхардал шалгах & Хангасан \\
\hline
\end{tabular}
\label{tab:functional_result}
\end{table}

\ \ \ \ \ \ \ \ \  \ \ \ \ Шалгалтын явцад систем эх сурвалжтай
баримт олдсон үед хариултын доор хууль болон холбоосыг харуулж байсан.
Харин өгөгдлийн санд шууд хамаарах мэдээлэл олдоогүй үед ``хангалттай
мэдээлэл олдсонгүй'' гэсэн fallback маягийн хариулт өгөх замаар
зохиомол мэдээлэл үүсгэхээс зайлсхийсэн. Энэ нь эрх зүйн чиглэлийн
чатботод hallucination багасгах зорилготой нийцэж байна.

\section{Хариултын чанарын ажиглалт}

\ \ \ \ \ \ \ \ \  \ \ \ \ Чанарын шалгалтаар хариултыг дараах
дөрвөн шалгуураар ажигласан: асуултад нийцсэн байдал, эх сурвалжийн
хамаарал, хариултын ойлгомжтой байдал, хийсвэр мэдээлэл үүсгэхээс
зайлсхийсэн байдал. Ерөнхий хууль зүйн нөхцөл асуусан үед систем
илүү тайлбарласан, хэрэглэгчид ойлгомжтой хариулт гаргаж байсан бол
яг зүйл заалтын дугаар асуух үед keyword search болон BM25-ийн нөлөө
илүү тод ажиглагдсан.

\begin{table}[H]
\renewcommand{\tablename}{Хүснэгт}
\centering
\caption{Хариултын чанарын ажиглалтын үр дүн}
\begin{tabular}{|l|l|}
\hline
\textbf{Шалгуур} & \textbf{Ажиглалт} \\
\hline
Асуултад нийцсэн байдал & Ерөнхий болон тодорхой асуултад тохирсон хариулт өгсөн \\
Эх сурвалжийн хамаарал & Хууль, URL, related laws/cases бүтэцтэй буцсан \\
Ойлгомжтой байдал & Хариулт хэрэглэгчид зориулсан тайлбар хэлбэртэй байсан \\
Hallucination бууруулах & Контекст сул үед fallback хариулт ашигласан \\
Follow-up ойлголт & Өмнөх ярианы түүхтэй асуултыг context болгон ашигласан \\
\hline
\end{tabular}
\label{tab:answer_quality}
\end{table}

\ \ \ \ \ \ \ \ \  \ \ \ \ Гэхдээ чанарын үнэлгээний энэ шат нь
голчлон гар ажиллагааны ажиглалт тул бүрэн тоон үнэлгээ гэж үзэх
боломжгүй. Цаашид 50-аас дээш асуулт-хариултын gold dataset бэлтгэж,
retrieval талд Recall@K, MRR, NDCG, generation талд эх сурвалжийн
үнэн зөв байдал болон хариултын бүрэн байдлыг оноогоор хэмжих
шаардлагатай.

\section{Туршилтын дүгнэлт}

\ \ \ \ \ \ \ \ \  \ \ \ \ Туршилтын үр дүнгээс харахад системийн
үндсэн архитектур ажиллаж, өгөгдөл боловсруулах pipeline, API service,
generation logic болон conversation persistence зэрэг гол хэсгүүд
шаардлагын түвшинд хэрэгжсэн байна. API талын автомат тестүүд бүрэн
амжилттай ажилласан нь online RAG pipeline-ийн үндсэн logic тогтвортой
байгааг харуулсан. Мөн локал ingestion туршилтаар 100 баримтыг алдаагүй
боловсруулж, ChromaDB-д индексжүүлсэн нь offline pipeline-ийн үндсэн
ажиллагааг баталгаажуулсан.

\ \ \ \ \ \ \ \ \  \ \ \ \ Үүний зэрэгцээ worker талын зарим автомат
тест унасан нь ingestion job-ийн тест орчин PostgreSQL service-ээс
хамааралтай хэвээр байгаа, мөн chunker-ийн тестүүд шинэ metadata
логиктой бүрэн нийцээгүй байгааг харуулж байна. Иймээс цаашдын
сайжруулалтад тестийн орчныг container-тэй уялдуулах, chunker-ийн
golden test dataset шинэчлэх, retrieval/generation чанарын тоон
үнэлгээний багц бүрдүүлэх ажлыг нэн тэргүүнд хийх шаардлагатай.

\chapter{Олж авсан сургамж}

\section{Монгол хэлний эрх зүйн өгөгдлийн онцлог}

\ \ \ \ \ \ \ \ \  \ \ \ \ Энэхүү ажлын явцад хамгийн түрүүнд ойлгосон зүйл нь хууль зүйн чатботын чанар зөвхөн хэлний загвараас бус, эх өгөгдлийн боловсруулалтын чанараас шууд хамаардаг явдал байв. Ялангуяа Монгол кирилл текстэд Unicode-ийн зөрүү, үл үзэгдэх тэмдэгт, эвдэрхий HTML бүтэц, зүйл заалтын дугаарлалтын олон янз хэлбэр зэрэг нь хайлтын чанарт шууд нөлөөлж байлаа. Иймээс parser, cleaner, normalizer, chunker гэсэн шат бүрийг тусад нь авч үзэж, Unicode NFC нормчилол, бүтцэд суурилсан хэсэгчлэл, metadata хавсаргалт зэрэг алхмуудыг сайтар шийдэх шаардлагатай болсон.

\ \ \ \ \ \ \ \ \  \ \ \ \ Энэ туршлагаас үзэхэд Монгол хэлний эрх зүйн системд retrieval-ийн чанарыг сайжруулах хамгийн үр дүнтэй арга нь илүү том загвар сонгохоос өмнө баримт бичгийн цэвэрлэгээ, логик бүтцийн задлал, citation-д хэрэгтэй metadata-г зөв бүрдүүлэх явдал юм. Өөрөөр хэлбэл асуултад зөв хариулах чадварын суурь нь чанартай context бүрдүүлэхэд оршиж байгааг энэхүү хэрэгжүүлэлт баталсан.

\section{Хосолсон хайлт ба туршилтад суурилсан тохируулга}

\ \ \ \ \ \ \ \ \  \ \ \ \ Хоёр дахь чухал сургамж нь хуулийн асуултад ганц төрлийн хайлт хангалтгүй байдаг явдал байв. Зарим асуултад зүйл, заалтын дугаар, хуулийн нэр, кейсийн дугаар зэрэг яг таарсан түлхүүр үг чухал болдог бол, заримд нь ерөнхий утга санааны ойролцоо байдал илүү чухал байлаа. Иймээс vector search болон PostgreSQL full-text search-ийг хамтад нь ашиглаж, үр дүнг RRF-ээр нэгтгэх хосолсон арга нь дан ганц semantic эсвэл дан ганц keyword хайлтаас илүү тогтвортой ажилласан.

\ \ \ \ \ \ \ \ \  \ \ \ \ Мөн retrieval-ийн дараах дахин эрэмбэлэлт, top-N context сонголт, score threshold, query rewrite, follow-up асуултын history ашиглалт зэрэг нь онолын таамаглалаар бус, бодит туршилт болон лог шинжилгээнд тулгуурлан тохируулагдах ёстойг олж харсан. Нэг тохиргоо бүх төрлийн асуултад адил сайн ажиллахгүй тул retrieval pipeline-ийг модульчлан зохион байгуулж, дараа дараагийн сайжруулалтад уян хатан байлгах нь зөв шийдэл болсон.

\section{Архитектурын салгалт ба найдвартай ажиллагаа}

\ \ \ \ \ \ \ \ \  \ \ \ \ Гурав дахь сургамж нь системийн найдвартай ажиллагаа зөвхөн API хариултын хурднаар хэмжигдэхгүй, харин ingestion pipeline, persistence, service хоорондын хамаарлыг хэрхэн зохион байгуулснаас ихээхэн шалтгаалдаг явдал байв. Worker, API, Web гэсэн салангид модулиудтай байснаар өгөгдөл цуглуулах, хайлт хийх, хэрэглэгчтэй харилцах урсгалууд хоорондоо хэт барьцалдахгүй болж, тус тусдаа сайжруулах боломж бүрдсэн. \texttt{packages/shared} дэх нэгтгэсэн schema, type-ууд нь frontend-backend интеграцийн алдааг эрс багасгасан.

\ \ \ \ \ \ \ \ \  \ \ \ \ Мөн offline pipeline дээр retry/backoff, checkpoint, dead-letter бүртгэл, chunk deduplication зэрэг механизмуудыг оруулснаар бодит орчинд удаан хугацаанд ажиллах суурь тавигдсан. Энэ нь хууль зүйн өгөгдөлтэй ажиллах системд зөвхөн ``хариулт гаргах'' чадвар бус, ``найдвартай шинэчлэгдэж, алдааг даван туулах'' чадвар адил чухал гэдгийг харуулсан. Иймээс цаашдын хөгжүүлэлтэд ч архитектурын салгалт, үйл ажиллагааны тогтвортой байдал хоёрыг гол зарчим хэвээр авч үзэх шаардлагатай.

%----------------------------------------------------------------------------------------
%   Дүгнэлт эндээс эхэлнэ
%----------------------------------------------------------------------------------------
\conclusion{Дүгнэлт}
\ \ \ \ \ \ \ \ \  \ \ \ \ Энэхүү дипломын ажлын хүрээнд Монгол хэл дээрх хууль зүйн асуултад эх сурвалжид тулгуурлан хариулах RAG суурьтай чатботын системийг хэрэгжүүлэв. Систем нь \texttt{legalinfo} болон \texttt{shuukh} эх сурвалжаас өгөгдөл цуглуулах worker pipeline, PostgreSQL ба вектор хадгалалтын давхарга, Fastify-д суурилсан retrieval ба generation API, мөн Next.js вэб хэрэглэгчийн интерфейс гэсэн үндсэн бүрэлдэхүүнтэйгээр ажиллаж байна. Хэрэгжүүлэлтийн түвшинд Монгол хэлний текст нормчилол, бүтэцт chunking, hybrid retrieval, heuristic reranking, conversation persistence зэрэг нь системийн бодит ажиллагааны гол тулгуур болсон.

\ \ \ \ \ \ \ \ \  \ \ \ \ Ажлын гол дүгнэлт нь хууль зүйн чатботын чанар зөвхөн LLM-ийн чадавхаас бус, харин өгөгдлийн боловсруулалт, хайлтын стратеги, context сонголт, эх сурвалжийн бүтэцжилтээс хамгийн их хамаардаг явдал юм. Ялангуяа Монгол хэлний онцлогтой эрх зүйн текст дээр Unicode нормчилол, metadata баяжуулалт, section-aware chunking, vector болон keyword хайлтын хослол нь хариултын чанарыг мэдэгдэхүйц сайжруулж байгааг энэхүү системийн хэрэгжүүлэлт харуулав.

\ \ \ \ \ \ \ \ \  \ \ \ \ Цаашид эх сурвалжийн хамрах хүрээг өргөтгөх, үнэлгээний тогтмол дата багц бүрдүүлэх, reranking болон generation шатны чанарын шалгуурыг улам боловсронгуй болгох, мөн хэрэглэгчийн туршлагыг сайжруулах нэмэлт боломжууд нээлттэй байна. Гэсэн хэдий ч одоогийн хэрэгжүүлэлт нь Монголын эрх зүйн мэдээллийн орчинд бодит эх сурвалжид суурилсан, өргөтгөх боломжтой, модульчлагдсан хууль зүйн чатботын суурь системийг амжилттай байгуулсан гэж дүгнэж болно.

%----------------------------------------------------------------------------------------
%   Дипломын номзүй, хавсралтын хэсэг эндээс эхэлнэ
%----------------------------------------------------------------------------------------

\singlespace
\addcontentsline{toc}{part}{НОМ ЗҮЙ}
\bibliographystyle{plain}
\bibliography{references}


%----------------------------------------------------------------------------------------
%   Хавсралтууд эндээс эхэлнэ
%----------------------------------------------------------------------------------------
\appendix
\addcontentsline{toc}{part}{ХАВСРАЛТ}

% Хавсралтын нэр. Хавсралт гэдэг үг агуулахгүй
\clearpage
\chapter{Кодын хэрэгжүүлэлт}
\markboth{}{}

\section{Хуваалцсан схем ба төрөл}

\ \ \ \ \ \ \ \ \  \ \ \ \ Системийн бүх модуль
\texttt{@legal-chatbot/shared} багцыг ашиглаж нэг ижил
өгөгдлийн гэрээг баримталдаг. Ингэснээр frontend, API,
worker талууд \texttt{conversationId}, \texttt{message},
\texttt{history}, \texttt{sources} зэрэг талбаруудыг өөр өөрөөр
тайлбарлах эрсдэл буурч, runtime болон compile-time түвшинд
зэрэг шалгагддаг.

\begin{lstlisting}[language=JavaScript,
caption={Shared багц дахь хүсэлтийн схем ба хариултын төрөл}]
// packages/shared/src/schemas/chat.schema.ts
import { z } from 'zod';

export const chatMessageSchema = z.object({
  role: z.enum(['user', 'assistant']),
  content: z.string().min(1, 'Message content cannot be empty'),
});

export const chatRequestSchema = z.object({
  conversationId: z.string().min(1).max(128).optional(),
  message: z.string().min(1).max(8000),
  history: z.array(chatMessageSchema).max(50).default([]),
});

// packages/shared/src/types/chat.types.ts
export interface ChatRequest {
  conversationId?: string;
  message: string;
  history: ChatMessage[];
}

export interface ChatResponse {
  conversationId: string;
  answer: string;
  sources?: Source[];
  relatedLaws?: RelatedLaw[];
  relatedCases?: RelatedCase[];
  confidence?: number;
}
\end{lstlisting}

\section{Монгол текст цэвэрлэгээ ба chunking}

\ \ \ \ \ \ \ \ \  \ \ \ \ \texttt{apps/worker} модульд
баримтын текстийг эхлээд Unicode түвшинд нормчилж, дараа нь
догол мөрөөр нь тасалж, давхцалтай chunk болгон хадгалдаг.
Энэ нь Монгол хэлний тэмдэгтийн эвдрэлээс үүсэх хайлтын
алдааг бууруулж, нэг зүйл заалтын төгсгөл дараагийн chunk-д
үргэлжилсэн тохиолдолд контекст алдагдахгүй байх нөхцөлийг
бүрдүүлсэн.

\begin{lstlisting}[language=JavaScript,
caption={Worker талын нормчлол ба баримт хэсэгчлэл}]
// apps/worker/src/processing/normalizer.ts
export function normalizeText(text: string): string {
  let normalized = text;
  normalized = normalizeUnicode(normalized);
  normalized = fixEncodingIssues(normalized);
  normalized = normalizePreservingParagraphs(normalized);
  normalized = normalizeMongolianText(normalized);
  return normalized;
}

// apps/worker/src/processing/chunker.ts
export function chunkDocument(doc, documentId, options = {}) {
  const paragraphs = doc.cleanedText
    .split(/\n\n+/)
    .filter((p) => p.trim().length > 0);

  const chunks = [];
  let currentText = '';

  for (const paragraph of paragraphs) {
    if (estimateTokens(currentText + ' ' + paragraph) > CHUNK_CONFIG.MAX_TOKENS
        && currentText.length > 0) {
      chunks.push(createChunk(currentText, chunks.length, documentId, 0, 0, metadata));
      currentText = getOverlapTail(currentText, CHUNK_CONFIG.OVERLAP_TOKENS)
        + ' ' + paragraph;
    } else {
      currentText = currentText ? currentText + '\n\n' + paragraph : paragraph;
    }
  }

  if (currentText.trim().length > 0) {
    chunks.push(createChunk(currentText, chunks.length, documentId, 0, 0, metadata));
  }
  return chunks;
}
\end{lstlisting}

\section{Embedding сервис}

\ \ \ \ \ \ \ \ \  \ \ \ \ Вектор үүсгэх хэсэг нь OpenAI болон
локал embedding гэсэн хоёр горимтой. Хэрэв \texttt{OPENAI\_API\_KEY}
өгөгдөөгүй бол систем автоматаар локал загвар руу шилждэг.
Мөн batch боловсруулах, retry, exponential backoff логик
оруулснаар олон chunk-тэй баримт индексжүүлэх үед тогтвортой
ажиллах нөхцөл бүрдсэн.

\begin{lstlisting}[language=JavaScript,
caption={Embedding service-ийн batching ба retry логик}]
// apps/worker/src/services/embedding.service.ts
export class EmbeddingService {
  constructor(config = {}) {
    const providerPreference = config.embeddingProvider ?? 'openai';
    const hasApiKey = Boolean(config.openaiApiKey);
    this.useLocal = providerPreference === 'local' || !hasApiKey;
    this.batchSize = config.batchSize ?? 100;
    this.maxRetries = config.maxRetries ?? 4;
  }

  async embedChunks(chunks) {
    const embedded = [];
    for (let i = 0; i < chunks.length; i += this.batchSize) {
      const batch = chunks.slice(i, i + this.batchSize);
      const vectors = await this.embedWithRetry(
        batch.map((c) => c.text),
        `batch-${Math.floor(i / this.batchSize) + 1}`,
      );
      embedded.push(...vectors);
    }
    return embedded;
  }

  private async embedWithRetry(texts, context) {
    let attempt = 0;
    while (attempt < this.maxRetries) {
      attempt += 1;
      try {
        return this.useLocal ? await this.embedLocal(texts)
                             : await this.embedOpenAI(texts);
      } catch (err) {
        if (attempt >= this.maxRetries || !this.isRetryableError(err)) throw err;
        await this.sleep(this.retryBaseMs * Math.pow(2, attempt - 1));
      }
    }
    throw new Error(`Embedding failed for ${context}`);
  }
}
\end{lstlisting}

\section{Hybrid хайлт ба дахин эрэмбэлэлт}

\ \ \ \ \ \ \ \ \  \ \ \ \ Хайлтын үндсэн логик нь
\texttt{apps/api/src/services/retrieval.service.ts} файлд
хэрэгжсэн. Асуулгыг эхлээд rewrite хийж, дараа нь vector хайлт,
keyword хайлт хоёрыг зэрэг ажиллуулна. Үр дүнг
\texttt{rrfService.fuse} функцээр нэгтгээд, хамгийн өндөр
магадлалтай хэсгүүдийг \texttt{rerankerService.rerank}
функцээр дахин эрэмбэлж generation шатанд дамжуулдаг.

\begin{lstlisting}[language=JavaScript,
caption={Hybrid retrieval: query rewrite, RRF fusion, rerank}]
// apps/api/src/services/retrieval.service.ts
export async function search(env, query, options = {}) {
  const baseQuery = extractPrimaryQuery(query);
  const rewrittenQuery = rewriteQuery(baseQuery);
  const searchQueries = buildSearchQueries(baseQuery, isShortQuery(baseQuery), rewrittenQuery);

  const embeddings = await Promise.all(
    searchQueries.map((q) =>
      embedQueryOpenAI(getOpenAIClient(env.OPENAI_API_KEY), env.OPENAI_EMBEDDING_MODEL, env.EMBEDDING_DIMENSION, q),
    ),
  );

  const vectorVariantResults = await Promise.all(
    embeddings.map((embedding) =>
      queryVectorStore(vectorEnv, embedding, OVERALL_TOP_K, { source: 'legalinfo' }),
    ),
  );
  const vectorResults = mergeVectorResults(vectorVariantResults.flat());

  let keywordResults = await keywordSearchService.search(baseQuery, 40, { source: 'legalinfo' });
  const rrfResults = rrfService.fuse(vectorResults, keywordResults);

  const fusedResults = rrfResults.map((r) => ({
    id: r.id,
    document: r.text,
    metadata: r.metadata,
    score: Math.max(r.vectorScore, r.keywordScore, Math.min(r.rrfScore * 10, 1)),
  }));

  const topCandidates = fusedResults.slice(0, RETRIEVAL_TOP_RESULTS);
  const rerankedTopChunks = await rerankerService.rerank(
    baseQuery,
    topCandidates,
    FINAL_TOP_RESULTS,
  );

  return buildStructuredRetrievalResult(rerankedTopChunks, fusedResults);
}
\end{lstlisting}

\begin{lstlisting}[language=JavaScript,
caption={RRF болон heuristic reranker-ийн үндсэн хэсэг}]
// apps/api/src/services/rrf.service.ts
export class RRFService {
  private k = 60;

  fuse(vectorResults, keywordResults) {
    const rrfMap = new Map();

    for (let i = 0; i < vectorResults.length; i++) {
      const rrfScore = 1 / (this.k + i + 1);
      rrfMap.set(vectorResults[i].id, { ...vectorResults[i], rrfScore, source: 'vector' });
    }

    for (let i = 0; i < keywordResults.length; i++) {
      const rrfScore = 1 / (this.k + i + 1);
      const existing = rrfMap.get(keywordResults[i].chunkId);
      if (existing) existing.rrfScore += rrfScore;
    }

    return Array.from(rrfMap.values()).sort((a, b) => b.rrfScore - a.rrfScore);
  }
}

// apps/api/src/services/reranker.service.ts
export class RerankerService {
  async rerank(query, chunks, topN = 8) {
    const queryTerms = query.toLowerCase().split(/\s+/).filter((term) => term.length > 1);
    const scored = chunks.map((chunk, originalRank) => {
      let score = chunk.score;
      const chunkText = `${String(chunk.metadata?.title ?? '')} ${chunk.document ?? ''}`.toLowerCase();
      const keywordMatches = queryTerms.reduce(
        (sum, term) => sum + (chunkText.match(new RegExp(term, 'g'))?.length || 0),
        0,
      );
      score += Math.min(keywordMatches * 0.03, 0.35);
      score += 0.05 * (1 - originalRank / chunks.length);
      return { ...chunk, score: Math.min(Math.max(score, 0), 1) };
    });
    return scored.sort((a, b) => b.score - a.score).slice(0, topN);
  }
}
\end{lstlisting}

\section{Chat API ба хариулт үүсгэх урсгал}

\ \ \ \ \ \ \ \ \  \ \ \ \ \texttt{registerChatRoute} функц нь
оролтын баталгаажуулалт, conversation ownership шалгалт,
retrieval plan үүсгэх, хайлт, rerank, generation, хадгалалт
гэсэн бүх алхмыг нэг урсгалд удирддаг. Одоогийн
хэрэгжүүлэлтэд follow-up асуултыг өмнөх түүхтэй холбох
логик тусгайлан орсон нь энгийн чат интерфейсийг бодит
харилцан ярианы систем болгоход чухал нөлөө үзүүлсэн.

\begin{lstlisting}[language=JavaScript,
caption={Fastify chat route-ийн үндсэн урсгал}]
// apps/api/src/routes/v1/chat.route.ts
export function registerChatRoute(app: FastifyInstance) {
  const handler = async (request, reply) => {
    const authUser = request.authUser;
    if (!authUser) {
      return reply.status(401).send({ error: 'UNAUTHORIZED', message: 'Authentication required.' });
    }

    const parseResult = chatRequestSchema.safeParse(request.body);
    if (!parseResult.success) {
      return reply.status(400).send({
        error: 'VALIDATION_ERROR',
        message: parseResult.error.errors.map((issue) => issue.message).join('; '),
      });
    }

    const { message } = parseResult.data;
    let history = parseResult.data.history;
    const conversationId = parseResult.data.conversationId ?? randomUUID();

    const retrievalPlan = buildRetrievalPlan(message, history);
    const retrievalResult = await search(app.env, retrievalPlan.query, {
      preferredLawIds: retrievalPlan.preferredLawIds,
    });

    const rerankedChunks = await rerankerService.rerank(
      retrievalPlan.query,
      retrievalResult.contextChunks,
      retrievalPlan.usesHistoryContext ? 10 : undefined,
    );

    const generationResult = await generate(
      app.env,
      message,
      rerankedChunks,
      retrievalPlan.relevantHistory,
      retrievalResult.relatedCases,
    );

    await addMessage(conversationId, 'user', message);
    await addMessage(conversationId, 'assistant', generationResult.answer, {
      confidence: generationResult.confidence,
      sources: retrievalResult.sources,
      relatedLaws: retrievalResult.relatedLaws,
      relatedCases: retrievalResult.relatedCases,
    });

    return reply.status(200).send({
      conversationId,
      answer: generationResult.answer,
      sources: retrievalResult.sources,
      relatedLaws: retrievalResult.relatedLaws,
      relatedCases: retrievalResult.relatedCases,
      confidence: generationResult.confidence,
    });
  };

  app.post('/v1/chat', { config: chatRateLimitConfig(), preHandler: app.authenticate }, handler);
}
\end{lstlisting}

\begin{lstlisting}[language=JavaScript,
caption={Generation service-ийн хамгаалалт ба fallback логик}]
// apps/api/src/services/generation.service.ts
export async function generate(env, query, contextChunks, history, relatedCases = []) {
  const primaryQuery = extractPrimaryQuery(query);
  const mode = detectQueryMode(primaryQuery);
  const resolvedIntent = resolveIntent(primaryQuery, history);

  if (resolvedIntent === 'unknown' && isGenericUnscopedLegalQuestion(primaryQuery)) {
    return {
      answer: NO_INFO_RESPONSE,
      confidence: 0,
      promptTokens: 0,
      completionTokens: 0,
      mode: 'no-info',
      suggestedQuestions: [],
    };
  }

  if (contextChunks.length === 0) {
    return buildFallbackResult(mode, resolvedIntent, primaryQuery);
  }

  const rerankedChunks = await rerankerService.rerank(primaryQuery, contextChunks, 8);
  if (rerankedChunks.length === 0) {
    return buildFallbackResult(mode, resolvedIntent, primaryQuery);
  }

  return buildLlmGroundedAnswer(env, primaryQuery, rerankedChunks, history, relatedCases);
}
\end{lstlisting}

\section{Frontend чат төлөвийн удирдлага}

\ \ \ \ \ \ \ \ \  \ \ \ \ Вэб талд \texttt{useChat} hook нь
API дуудах, conversation сонгох, мессежүүдийг state-д хадгалах,
алдаа болон retry логикийг төвлөрүүлсэн. Ингэснээр UI
компонентууд зөвхөн дүрслэлдээ төвлөрч, өгөгдлийн урсгал нь
нэг цэгээс удирдагдах болсон.

\begin{lstlisting}[language=JavaScript,
caption={Next.js талын chat state ба message dispatch}]
// apps/web/lib/useChat.ts
export function useChat() {
  const [conversationId, setConversationId] = useState(null);
  const [messages, setMessages] = useState([]);
  const [isLoading, setIsLoading] = useState(false);
  const [error, setError] = useState(null);

  const sendMessage = useCallback(async (content) => {
    const trimmed = content.trim();
    if (!trimmed) return;

    const history = messagesRef.current.map((message) => ({
      role: message.role,
      content: message.content,
    }));

    setMessages((current) => [
      ...current,
      { id: uid(), role: 'user', content: trimmed, timestamp: Date.now() },
    ]);
    setIsLoading(true);

    try {
      const response = await sendChatMessage({
        conversationId: conversationIdRef.current ?? undefined,
        message: trimmed,
        history,
      }, accessTokenRef.current);

      setConversationId(response.conversationId);
      setMessages((current) => [
        ...current,
        {
          id: uid(),
          role: 'assistant',
          content: response.answer,
          sources: response.sources,
          relatedCases: response.relatedCases,
          relatedLaws: response.relatedLaws,
          confidence: response.confidence,
          timestamp: Date.now(),
        },
      ]);
    } catch (reason) {
      setError(reason instanceof Error ? reason.message : 'Алдаа гарлаа.');
    } finally {
      setIsLoading(false);
    }
  }, []);
}
\end{lstlisting}

\end{document}
